\documentclass[11pt,reqno]{amsart}

\usepackage[letterpaper,margin=1in]{geometry}
\usepackage[T1]{fontenc}
\usepackage{newtxtext,newtxmath}
\usepackage{microtype}
\usepackage{amsmath,mathtools}
\usepackage[authoryear,round]{natbib}
\usepackage{enumitem}
\usepackage{xurl}
\usepackage[hidelinks]{hyperref}
\usepackage{bookmark}
\usepackage{fancyhdr}

\setcounter{secnumdepth}{2}
\setcounter{tocdepth}{2}
\setlength{\parindent}{1.45em}
\setlength{\parskip}{0pt}
\setlist[itemize]{leftmargin=2.05em,itemsep=0.12em,topsep=0.32em}
\setlist[enumerate]{leftmargin=2.15em,itemsep=0.16em,topsep=0.36em}
\setlength{\emergencystretch}{3em}
\allowdisplaybreaks
\providecommand{\tightlist}{\setlength{\itemsep}{0pt}\setlength{\parskip}{0pt}}

% Use unobtrusive outside folios without running heads: lower left on even
% pages and lower right on odd pages.
\pagestyle{fancy}
\fancyhf{}
\fancyfoot[LE,RO]{\thepage}
\renewcommand{\headrulewidth}{0pt}
\renewcommand{\footrulewidth}{0pt}
\setlength{\footskip}{24pt}

\newcommand{\Prob}[1]{\mathbb{P}\!\left(#1\right)}
\newcommand{\Exp}[1]{\mathbb{E}\!\left[#1\right]}
\newcommand{\gap}{\chi(G_n)-\zeta(G_n)}
\newcommand{\displayheading}[1]{\par\addvspace{0.68\baselineskip}\noindent{\scshape #1.}\par\nobreak\smallskip\noindent}

\theoremstyle{plain}
\newtheorem*{maintheorem}{Main theorem}
\newtheorem{theorem}{Theorem}[section]
\newtheorem{proposition}[theorem]{Proposition}
\newtheorem{lemma}[theorem]{Lemma}
\newtheorem{corollary}[theorem]{Corollary}

\theoremstyle{remark}
\newtheorem{remark}[theorem]{Remark}

\hypersetup{
  pdftitle={A Full-Sequence Quantitative Gap Between the Chromatic and Cochromatic Numbers of a Random Graph},
  pdfauthor={Samuil Petkov},
  pdfcreator={LaTeX}}

\title[Chromatic and cochromatic numbers]{A Full-Sequence Quantitative Gap
Between the Chromatic and Cochromatic Numbers of a Random Graph}
\author[Samuil Petkov]{Samuil Petkov}
\address{\'Ecole normale sup\'erieure, Universit\'e PSL, Paris, France}
\email{samuil.petkov@ens.psl.eu}
\subjclass[2020]{Primary 05C80; Secondary 05C15, 60C05}
\keywords{random graph, chromatic number, cochromatic number, second moment method, configuration model}

\begin{document}



\begin{abstract}
Let $G_n\sim G(n,1/2)$. We prove, along the full sequence of integers $n$, that
\[
  \chi(G_n)-\zeta(G_n)=\Omega\!\left(\frac{n}{(\log n)^3}\right)
\]
with probability tending to one, where $\chi$ and $\zeta$ denote the chromatic
and cochromatic numbers. This resolves the question of Erd\H{o}s and Gimbel
asking whether this difference tends to infinity with high probability. The
main difficulty is to obtain estimates that remain uniform when the natural
class-size cutoff changes. We use a signed profile supported on four
consecutive class sizes, an exact identity for the overlap of two such
profiles, and a matching restriction that controls the remaining cycle-space
sum. A bounded-differences argument then amplifies the resulting second-moment
witness to a high-probability cocoloring.
\end{abstract}

\maketitle
\thispagestyle{fancy}

\section*{Introduction}
\addcontentsline{toc}{section}{Introduction}
\label{sec:introduction-v3}

A \emph{cocoloring} of a graph $G$ is a partition of $V(G)$ into nonempty
classes, each inducing either an edgeless graph or a complete graph. The least
number of classes in such a partition is the \emph{cochromatic number}
$\zeta(G)$. Since an ordinary proper coloring is a cocoloring using only
independent-set classes, one always has $\zeta(G)\le\chi(G)$.

Write $[n]:=\{1,\ldots,n\}$, and let $G_n\sim G(n,1/2)$ be the labeled random
graph on $[n]$. All logarithms are natural unless a base is displayed. Erd\H{o}s and
Gimbel asked whether
\[
  \chi(G_n)-\zeta(G_n)\longrightarrow\infty
\]
with high probability \citep[p.~263]{erdos-gimbel-1993}. The question was
restated by Gimbel \citep[Section~7.4]{gimbel-2016} and is cataloged as
Erd\H{o}s Problem~625 \citep{bloom-erdos625}. We obtain the following
quantitative full-sequence statement.

\begin{maintheorem}
For $G_n\sim G(n,1/2)$,
\[
  \mathbb P\!\left(
    \chi(G_n)-\zeta(G_n)
    \ge
    \frac{(\log 2)^2}{32}
    \log\!\left(\frac{200}{153}\right)
    \frac{n}{(\log n)^3}
  \right)
  \longrightarrow 1.
\]
In particular, $\chi(G_n)-\zeta(G_n)$ tends to infinity with high probability
along the full sequence of integers $n$.
\end{maintheorem}

The phrase \emph{full sequence} is substantive. The natural class-size cutoff
jumps whenever the independence-number phase crosses an integer. At such a
jump, changing one admissible class size changes the optimizing profile and
its first-moment root. An estimate proved only for a density-one set of phase
values does not automatically extend to every integer $n$. All asymptotic
estimates in this paper are therefore uniform up to both endpoints of every
phase interval.

Unless stated otherwise, every $o(1)$ below is a deterministic error uniform
over the complete phase interval. We keep the high-skeleton and residual
attachment errors separate until the final second-moment ledger, which makes
the uniformity and the absence of double counting explicit.

\subsection*{Relation to previous work}

The first-order asymptotic for the chromatic number of a dense random graph was
proved by Bollob\'as \citep{bollobas-1988}, following the early work of
Grimmett and McDiarmid \citep{grimmett-mcdiarmid-1975}; later refinements
include \citet{mcdiarmid-1990}, \citet{panagiotou-steger-2009}, and
\citet{heckel-2018}. The cochromatic number belongs to the theory of generalized
chromatic numbers associated with hereditary graph properties
\citep{scheinerman-1992,bollobas-thomason-1995}.

For the difference $\chi(G_n)-\zeta(G_n)$, Heckel and, independently, Steiner
related divergence to the nonconcentration of the chromatic number
\citep{heckel-2024-question,steiner-2024}. Heckel subsequently proved a
near-linear lower bound for a phase-dependent set containing approximately
$95\%$ of the integers \citep{heckel-2025-difference}. The unresolved part is
precisely the exceptional phase. We retain the signed first-moment idea and
rare-seed amplification from that work, but replace the pairwise sign estimate
by an exact overlap identity. One four-size profile covers every phase; exact
endpoint-table summation and residual cycle-space restriction make its second
moment uniform.

\subsection*{Contributions and main ideas}

The proof has four reusable components.

\begin{enumerate}
\item \textbf{Uniform root separation.} We compare the ordinary coloring root
with the signed cocoloring root supported on four consecutive class sizes. The
finite support comparison is exact before passing to the limiting Gaussian
dual. We then compare the finite target with the phase-center target through
one explicit, phase-uniform error. The resulting separation has order
$n/(\log n)^3$.

\item \textbf{Exact sign compatibility.} For two signed witnesses, the sum over
all compatible sign assignments factors into local cell rewards and the size
of a binary cycle space. No sign correlation is estimated before this exact
identity has been extracted.

\item \textbf{Large cells without chosen completions.} Cells above half the
class-size cap form a matching. We sum the entire physical matching fiber,
compare it with full containment, and apply the ambient falling-factorial
bound once to the total deficit. The resulting reference weights are regrouped
exactly by their endpoint tables.

\item \textbf{Residual restriction and amplification.} After the large cells
have been exposed, deletion outside their matching is injective on even
residual edge sets. This converts the residual cycle-space sum into a product
bound. The high-skeleton and residual logarithmic costs are kept separate and
added only after the exact conditional decomposition. Paley--Zygmund supplies
a rare signed witness, and a one-sided bounded-differences argument amplifies
it to a high-probability cocoloring.
\end{enumerate}

The four consecutive sizes used here form a small fixed support that spans the
entire phase interval and leaves a uniform entropy advantage over ordinary
coloring. Here $\delta\in[0,1]$ is the rounding phase and $D_4(\delta)$ is the
finite-support entropy loss defined in Section~5. Lemma~\ref{lem:entropy-gap}
proves the uniform certificate
\[
  \log 2-D_4(\delta)
  >
  \log\!\left(\frac{200}{153}\right).
\]
The later factor $1/8$ leaves a factor-four margin over the displayed
coefficient $1/32$. Thus
the finite support is not merely a technical truncation: it is the device that
converts phase dependence into a full-sequence theorem.

\subsection*{Guide to the argument}

The \emph{location step} identifies two continuous first-moment roots. The
signed witness is placed at a tangent-rounded midpoint, retaining half of their
leading separation while satisfying the exact profile conservation laws. The
phase, finite-dual, target-displacement, and slope errors are uniform in the
complete phase.

The \emph{existence step} studies the normalized second moment of this one
signed profile. Partial diagonals are handled first. The remaining overlap
tables are divided into canonical high cells and residual cells. Sections~8
and~9 keep these two accounts disjoint: no residual factor is paid in the
high-cell estimate, and no high-cell multiplicity is paid again in the
residual estimate. The logarithm of the final upper bound is explicitly
\[
  \Lambda_n
  =\Gamma_n^{\mathrm{skel}}+\Gamma_n^{\mathrm{att}}
  =o\!\left(\frac{n}{(\log n)^4}\right).
\]

The \emph{amplification step} converts this second-moment seed into a typical
statement. The largest induced subgraph admitting the prescribed number of
cocoloring classes is one-Lipschitz under vertex-block exposure. Its
concentration, together with a simultaneous coloring bound for every leftover
set, makes the amplification uniform in the seed exponent.

Section~1 fixes the phase notation. Sections~2--5 locate and compare the two
roots. Section~6 proves the exact signed-overlap identity, and Section~7 treats
partial diagonals. Sections~8 and~9 contain the large-cell and residual
estimates and the explicit global error ledger. Section~10 proves the
amplifier. The final section combines the amplified cocoloring event with the
ordinary chromatic lower-tail event and applies the uniform entropy
certificate. A concise note after the proof identifies the exact
machine-checked theorem and its build.

\section*{Notation and proof objects}
\addcontentsline{toc}{section}{Notation and proof objects}
\label{sec:conventions-v3}

We collect here only the conventions that are used throughout several later
sections. More local notation is introduced where it first becomes necessary.
This keeps the main argument self-contained without forcing the reader to
retain a large dictionary before the proof begins.

\paragraph{Asymptotic conventions.}
An event holds \emph{with high probability} if its
probability tends to one as $n\to\infty$. Whenever a phase parameter is
present, $o(1)$ denotes one deterministic error sequence that is uniform over
the full phase interval, including sequences approaching either endpoint.
For integers $x,r\ge0$, we use the falling factorial
\[
  (x)_r=x(x-1)\cdots(x-r+1),
  \qquad
  (x)_0=1,
\]
and set $(x)_r=0$ for $r>x$. A quotient involving falling factorials is used
only after the denominator has been shown to be nonzero.

\paragraph{Signed witnesses and profiles.}
A \emph{signed cocoloring witness} is a partition together with an $I$- or
$K$-mark on each class. An $I$-marked class must be independent and a
$K$-marked class must be complete. The marks are counting data; they do not
define a new graph invariant. In the four-size profile every class has size at
least two for sufficiently large $n$, so a realized class cannot satisfy both
requirements. Forgetting the marks therefore recovers its cocoloring without
multiplicity.

A \emph{profile} is a finitely supported sequence $(k_s)_{s\ge1}$ of
nonnegative integers, where $k_s$ is the number of classes of size $s$. It is
feasible on $n$ vertices when
\[
  \sum_s k_s=k,
  \qquad
  \sum_s s k_s=n.
\]
The ordinary first moment counts partitions into independent classes. The
signed first moment counts the same partitions together with one $I/K$ mark
per class.

\paragraph{Overlap table and support graph.}
Let $(A_a)_a$ and $(B_b)_b$ be two ordered profile partitions. Their overlap
table is
\[
  r_{ab}=|A_a\cap B_b|.
\]
Its row and column sums are the class sizes of the two partitions. It may also
be viewed as the cell-count table of a bipartite configuration model: every
underlying vertex pairs one labeled row stub with one labeled column stub.

The \emph{support graph} $H(r)$ is the simple bipartite graph whose vertices
are the row and column classes incident to an edge and whose edge set is
\[
  E(H(r))=\{(a,b):r_{ab}\ge2\}.
\]
A subset of its edges is \emph{even} if every vertex has even degree. Writing
$c(H)$ for the number of connected components, we put
\[
  \beta(H)=|E(H)|-|V(H)|+c(H).
\]
Then $\beta(H)$ is the dimension of the binary cycle space, so the number of
even edge sets is $2^{\beta(H)}$. The threshold two is exact: a cell of size
zero or one contains no edge internal to both partition classes, whereas a
cell of size at least two forces the marks on its row and column classes to
agree.

\paragraph{Three matching objects.}
The proof uses three distinct objects that should not be conflated.
\begin{enumerate}
\item The \emph{configuration matching} is the perfect matching of all labeled
row and column stubs induced by the common vertex set.
\item A \emph{physical partial matching} is a set of selected stub pairs inside
specified overlap cells.
\item A \emph{block support} is a matching between row-class slots and
column-class slots; its edges indicate the cells declared canonical and high.
\end{enumerate}
The product formula for physical cell counts relies on the third object being
a matching. Distinct selected cells then use disjoint row and column stub
families.

\paragraph{High cells, deficits, and endpoint tables.}
Let $U$ be the largest class size in the four-size profile and put
$R_0=\lfloor U/2\rfloor$. A cell is \emph{high} when its multiplicity exceeds
$R_0$. Two high cells cannot share a row or a column, so the high cells form a
canonical block support. For a selected cell $e$, let $s_e,t_e$ be its endpoint
sizes and let $j_e$ be its realized multiplicity. Set
\[
  m_e=\min\{s_e,t_e\},
  \qquad
  h_e=m_e-j_e.
\]
Here $m_e$ is the full-containment multiplicity, $j_e$ is the realized
multiplicity, and $h_e$ is the deficit. Highness implies $2h_e<m_e$.

Section~8 uses the bijection $i\mapsto i+2$ from $\{0,1,2,3\}$ to the deficit
labels $\{2,3,4,5\}$. For $i\in\{0,1,2,3\}$, define
$u_i:=\alpha-(i+2)$ and $\kappa_i:=k_{i+2}$.
Thus $u_i$ is the endpoint size and $\kappa_i$ is the number of available
blocks corresponding to deficit label $i+2$. For a block support $P$, the
endpoint table $L(P)=(L_{ij})$ records how many support edges join types $i$
and $j$. A table is \emph{feasible} when its type-$i$ row sum and type-$j$
column sum are at most $\kappa_i$ and $\kappa_j$, respectively. Every block
support has a feasible table, and every feasible table is attained by a
block-slot matching.

\paragraph{Separation of accounts.}
The \emph{bare high-skeleton weight} contains the incidence probability and
local rewards of the exposed high cells. It contains neither the rewards of
unexposed cells nor the residual binary cycle-space factor. Conditional on a
high skeleton, the expectation of those remaining factors is its
\emph{residual attachment}. Sections~8 and~9 preserve this separation exactly:
a factor paid in one account is never paid again in the other.

Every load-bearing estimate follows the same order. We first state the finite
identity, then apply deterministic inequalities, then take phase asymptotics,
and only at the end deduce a probability statement. This order is useful both
for mathematical auditing and for the theorem-by-theorem formalization.

\section{Phase notation and elementary estimates}\label{notation-and-elementary-facts}

For an integer \(s\), let

\[
 \mu_s(v)=\binom vs2^{-\binom s2},\qquad \mu_s=\mu_s(n). \tag{1.1}
\]

We use the following standard inequalities in their stated forms.

\begin{enumerate}
\def\labelenumi{\arabic{enumi}.}
\tightlist
\item
  For integers \(m\ge 1\), \[
   \log(m!)=m\log m-m+\frac12\log(2\pi m)+O(1/m).          \tag{1.2}
  \] Consequently, with \(0\log 0=0\), one has
  \(\log(m!)=m\log m-m+O(\log(m+1))\) uniformly for
  \(m\in\mathbb N_0\), with an absolute implied constant.
\item
  If a random variable \(Y\) is a function of \(r\)
  independent blocks and changing one block changes \(Y\) by at
  most one, then \[
   \Prob{\lvert Y-\Exp{Y}\rvert\ge t}
   \le2\exp(-2t^2/r).                                    \tag{1.3}
  \] We will use the corresponding one-sided bounds as well.
\item
  If \(Z\ge 0\) and \(0<\Exp{Z}<\infty\), then \[
   \Prob{Z>0}\ge\frac{\Exp{Z}^2}{\Exp{Z^2}}.             \tag{1.4}
  \]
\item
  If \(X\sim\operatorname{Bin}(m,1/2)\), then \[
   \Prob{X\le m/4}\le e^{-m/16}.                         \tag{1.5}
  \] Indeed, exponential Markov with \(t=\log 3\) bounds the
  probability by
  \([3^{1/4}(2/3)]^m\le e^{-m/16}\).
\item
  For every nonnegative random variable \(X\) and \(a>0\),
  \(\Prob{X\ge a}\le \Exp{X}/a\).
\end{enumerate}

The first is Stirling's estimate, the second is McDiarmid's bounded-
differences inequality, and the third is the zero-threshold case of
Paley--Zygmund. The fourth is the only binomial-tail estimate used
below; the fifth is Markov's inequality. The bounded-differences
formulation is the one recorded by
\citet[Theorem~3.1]{mcdiarmid-1989}.



\section{The complete independence-number phase}
\label{the-complete-independence-number-phase}

Put
\[
  q:=\log 2,
  \qquad L:=\log n,
  \qquad \ell:=\log\log n,
  \qquad C:=1+\log q-q.
\]
The usual independence-number center can then be written without a change of
logarithm base as
\begin{equation*}
  \alpha_0
  =2\log_2 n-2\log_2\log_2 n+2\log_2(\mathrm e/2)+1
  =\frac{2(L-\ell+C)}q+1.
  \tag{2.1}
\end{equation*}
Let
\[
  \alpha:=\lfloor\alpha_0\rfloor,
  \qquad
  \delta:=\alpha_0-\alpha,
  \qquad
  b:=1-\delta.
\]
Thus $0\le\delta<1$, $0<b\le1$, and
\[
  \alpha=\frac{2(L-\ell+C)}q+b.
\]
When emphasizing dependence on $n$, we write $\delta_n:=\delta$.
All estimates below are uniform for the closed parameter range
$0\le\delta\le1$; this includes integer sequences approaching either endpoint
of the actual half-open phase interval.

\begin{lemma}[Uniform phase expansion and adjacent-size control]
There is a bounded continuous function $K\colon[0,1]\to\mathbb R$ and an
absolute constant $C_{\mathrm{ph}}>0$ such that, for all sufficiently large
$n$, the quantity
\[
  \varepsilon_n^{\mathrm{ph}}
  :=C_{\mathrm{ph}}\frac{1+\ell^2}{L}
  \longrightarrow0
\]
satisfies the following expansion at the induced phase
$\delta=\alpha_0-\lfloor\alpha_0\rfloor$:
\begin{equation*}
  \log\mu_\alpha
  =
  \delta L
  +\left(\frac2q-\frac12-\delta\right)\ell
  +K(\delta)
  +E_n(\delta),
  \qquad
  |E_n(\delta)|\le\varepsilon_n^{\mathrm{ph}}.
  \tag{2.2}
\end{equation*}
The error bound is uniform over all induced phase values, and therefore also
along sequences approaching either endpoint of $[0,1]$. Moreover, uniformly
over those phase values,
\begin{equation*}
\begin{split}
  \log\mu_{\alpha+2}
  &=(\delta-2)L
    +\left(\frac2q+\frac32-\delta\right)\ell+O(1),\\
  \mu_{\alpha+2}
  &\le \exp\{-L+C_2\ell\}=o(1)
\end{split}
  \tag{2.3}
\end{equation*}
for an absolute $C_2$, while
\begin{equation*}
  \mu_{\alpha-2}
  \ge c\,n^2(\log n)^{2/q-5/2}
  \tag{2.4}
\end{equation*}
for an absolute $c>0$ and all sufficiently large $n$.
\end{lemma}

\begin{proof}
Because $\alpha=O(L)$ uniformly in the phase, one has $\alpha<n/2$ for all
sufficiently large $n$. For $0\le x\le1/2$,
$-2x\le\log(1-x)\le-x$. Hence
\[
  \log(n)_\alpha
  =\alpha L+\sum_{j=0}^{\alpha-1}\log\!\left(1-\frac jn\right)
  =\alpha L+R_{n,\alpha},
  \qquad
  |R_{n,\alpha}|\le\frac{\alpha^2}{n}.
\]
Using the uniform Stirling remainder
\[
  \log(\alpha!)
  =\alpha\log\alpha-\alpha
   +\frac12\log(2\pi\alpha)+O(1/\alpha)
\]
in
\[
  \log\mu_\alpha
  =\log(n)_\alpha-\log(\alpha!)-q\binom\alpha2
\]
gives
\begin{equation*}
  \log\mu_\alpha
  =
  \alpha\left(
    L-\log\alpha+1-\frac q2(\alpha-1)
  \right)
  -\frac12\log(2\pi\alpha)
  +O(1/L).
  \tag{2.5}
\end{equation*}
All implied constants and remainder bounds here are uniform in $\delta$.

Write
\[
  \alpha=\frac{2L}{q}(1+x_n),
  \qquad
  x_n:=\frac{-\ell+C+qb/2}{L}.
\]
Uniformly for the actual range $0<b\le1$ (and on its closure), one has
$|x_n|\le1/2$ for all sufficiently large $n$. The Taylor formula
$\log(1+x)=x+O(x^2)$ therefore gives
\[
  \log\alpha
  =\ell+q-\log q
   +\frac{-\ell+C+qb/2}{L}
   +O\!\left(\frac{\ell^2}{L^2}\right).
\]
Since
\[
  \frac q2(\alpha-1)=L-\ell+C-\frac{q\delta}{2},
\]
we obtain the uniform expansion
\begin{equation*}
  L-\log\alpha+1-\frac q2(\alpha-1)
  =
  \frac{q\delta}{2}
  +\frac{\ell-C-qb/2}{L}
  +O\!\left(\frac{\ell^2}{L^2}\right).
  \tag{2.6}
\end{equation*}
Multiplying by
$\alpha=2(L-\ell+C)/q+b$ and collecting terms yields
\[
\begin{split}
  \alpha\left(
    L-\log\alpha+1-\frac q2(\alpha-1)
  \right)
  ={}&
  \delta L+\left(\frac2q-\delta\right)\ell\\
  &+\delta C+\frac q2\delta(1-\delta)
    -\frac{2C}{q}-(1-\delta)
    +O\!\left(\frac{1+\ell^2}{L}\right).
\end{split}
\]
Also
\[
  -\frac12\log(2\pi\alpha)
  =-\frac12\ell-\frac12\log(2\pi)-\frac q2
    +\frac12\log q+O(\ell/L).
\]
Thus (2.2) holds with
\begin{equation*}
\begin{split}
  K(\delta):={}&
  \delta C+\frac q2\delta(1-\delta)-\frac{2C}{q}-(1-\delta)\\
  &-\frac12\log(2\pi)-\frac q2+\frac12\log q.
\end{split}
  \tag{2.7}
\end{equation*}
This function is continuous, hence bounded, on $[0,1]$. The combined remainder
in (2.5)--(2.7) has absolute value at most
$C'(1+\ell^2)/L$ for an absolute $C'$. Taking
$C_{\mathrm{ph}}\ge C'$ gives
$|E_n(\delta)|\le\varepsilon_n^{\mathrm{ph}}$ in (2.2).

For integers $1\le s<n$, the adjacent-size ratios are exact:
\begin{equation*}
  \frac{\mu_{s+1}}{\mu_s}
  =\frac{n-s}{s+1}2^{-s},
  \qquad
  \frac{\mu_{s-1}}{\mu_s}
  =\frac{s}{n-s+1}2^{s-1}.
  \tag{2.8}
\end{equation*}
The displayed formula for $\alpha$ gives the exact phase representation
\[
  2^\alpha
  =\exp(q\alpha)
  =\exp(2C+qb)\frac{n^2}{L^2}.
\]
Since $0<b\le1$, the multiplicative factor $\exp(2C+qb)$ is bounded
above and below by positive absolute constants. Therefore, for every fixed
$M$, uniformly over integers $s$ with $|s-\alpha|\le M$,
\[
  \frac{\mu_{s+1}}{\mu_s}=\Theta(L/n),
  \qquad
  \frac{\mu_{s-1}}{\mu_s}=\Theta(n/L).
\]
In particular,
\begin{equation*}
  \mu_{\alpha+2}
  =\mu_\alpha\,\Theta(L^2/n^2),
  \qquad
  \mu_{\alpha-2}
  =\mu_\alpha\,\Theta(n^2/L^2).
  \tag{2.8a}
\end{equation*}
Taking logarithms and using (2.2) proves the first line of (2.3). Since
$\delta\le1$, its leading term is at most $-L$, and the coefficient of
$\ell$ is uniformly bounded; this proves the second line.

For the lower adjacent size, (2.2) and (2.8a) give
\[
\begin{split}
  \log\mu_{\alpha-2}
  ={}&2L+\left(\frac2q-\frac52\right)\ell
      +\delta(L-\ell)+K(\delta)+O(1).
\end{split}
\]
For large $n$, $L-\ell>0$, while $K$ is bounded below. Exponentiating gives
(2.4), with one absolute constant $c$ and one eventuality threshold valid for
the complete phase.
\end{proof}

Let $X_{\alpha+2}$ denote the number of independent sets of size
$\alpha+2$. Then
$\mathbb E X_{\alpha+2}=\mu_{\alpha+2}$, and the event
$\alpha(G_n)>\alpha+1$ implies $X_{\alpha+2}\ge1$. Consequently Markov's
inequality and (2.3) give the deterministic cap error
\begin{equation*}
  \mathbb P\bigl(\alpha(G_n)>\alpha+1\bigr)
  \le\mu_{\alpha+2}
  =:\varepsilon_n^{\mathrm{cap}}
  \longrightarrow0
  \tag{2.9}
\end{equation*}
uniformly along the full sequence of integers.


\section{Continuous profile roots}
\label{continuous-profile-roots}

For an integer class size $u\ge1$, put
\[
  d_u:=2^{\binom u2}u!.
\]
Write $u=\alpha-i$, so $i$ is the deficit from the phase center $\alpha$.
We use
\begin{equation*}
  S_+:=\{-1,0,1,2,\ldots\},
  \qquad
  S_4:=\{2,3,4,5\}.
  \tag{3.1}
\end{equation*}
At finite $n$, the unrestricted support is
$S_+^{(n)}:=\{-1,0,\ldots,\alpha-1\}$: the lower endpoint comes from the
cap event (2.9), and the upper endpoint from positivity of class sizes.
Throughout every finite-$n$ profile, maximization, or sum indexed by $S_+$,
this cutoff is understood.

For $i\in S_+^{(n)}$, define the curved score
\begin{equation*}
  h_n(i):=
  \begin{cases}
    -\dfrac{\log 2}{2}i^2
      +\displaystyle\sum_{r=0}^{i-1}\log\!\left(1-\dfrac r\alpha\right),
      & i\ge0,\\[1ex]
    -\dfrac{\log 2}{2}+\log\!\left(\dfrac{\alpha}{\alpha+1}\right),
      & i=-1.
  \end{cases}
  \tag{3.1a}
\end{equation*}

For a support $S\in\{S_+^{(n)},S_4\}$ and a nonnegative real profile
$\mathbf k=(k_i)_{i\in S}$, put $k:=\sum_{i\in S}k_i$. When $k>0$, define
\begin{equation*}
  L_{\mathbf k}
  :=n\log n-n+k-\sum_{i\in S} k_i\log(k_i d_{\alpha-i}),
  \tag{3.2}
\end{equation*}
with the convention $0\log0=0$, and let $L_S(n,k)$ be its maximum over
real $k_i\ge0$ satisfying
\begin{equation*}
  \sum_{i\in S}k_i=k,
  \qquad
  \sum_{i\in S}(\alpha-i)k_i=n.
  \tag{3.3}
\end{equation*}
For $k$ feasible for the support $S$ under consideration, set
\begin{equation*}
  s:=\frac nk,
  \qquad
  T:=\alpha-s=\alpha-\frac nk.
  \tag{3.4}
\end{equation*}
With $p_i:=k_i/k$, the constraints give $T=\sum_{i\in S}ip_i$; thus $T$ is
the mean deficit of the normalized profile.

\begin{lemma}[Uniform root, slope, and finite-dual package]
Let
\begin{equation*}
  s_0:=\alpha_0-1-\frac2{\log 2},
  \qquad
  T_0:=\alpha-s_0
      =1+\frac2{\log 2}-\delta.
  \tag{3.5}
\end{equation*}
There is a fixed $A_*>0$ such that, for
$S\in\{S_+,S_4\}$ and $0\le c\le\log 2$, the equation
\[
  L_S(n,k)+ck=0
\]
has, for all sufficiently large $n$, a unique zero $r_{S,c}$ in the
  corridor
\[
  \left|\frac n{r_{S,c}}-s_0\right|
  \le A_*\frac{\log\log n}{\log n}.
\]
Uniformly in $S$, $c$, and the phase,
\begin{equation*}
  \frac n{r_{S,c}}
  =s_0+O\!\left(\frac{\log\log n}{\log n}\right),
  \qquad
  \alpha-\frac n{r_{S,c}}
  =T_0+O\!\left(\frac{\log\log n}{\log n}\right).
  \tag{3.6}
\end{equation*}

For every fixed $A>0$, uniformly over $S\in\{S_+,S_4\}$,
$0\le c\le\log 2$, and $k$ feasible for $S$ with
$|n/k-s_0|\le A\log\log n/\log n$,
\begin{equation*}
  \frac{\partial}{\partial k}\{L_S(n,k)+ck\}
  =\frac2{\log 2}(\log n)^2
   +O_A((\log n)(\log\log n)).
  \tag{3.7}
\end{equation*}
In particular, the deterministic normalized slope error
\[
  \varepsilon_{n,A}^{\mathrm{slope}}
  :=
  \max_{\substack{S\in\{S_+,S_4\}\\0\le c\le\log 2}}
  \sup_{\substack{k\ \mathrm{feasible\ for}\ S\\
                   |n/k-s_0|\le A\log\log n/\log n}}
  \left|
    \frac{1}{(\log n)^2}
    \frac{\partial}{\partial k}\{L_S(n,k)+ck\}
    -\frac2{\log 2}
  \right|
\]
satisfies $\varepsilon_{n,A}^{\mathrm{slope}}\to0$.

For two supports $R,S\in\{S_+,S_4\}$ at a $k$ feasible for both supports,
\begin{equation*}
  \frac{L_R(n,k)-L_S(n,k)}{k}
  =\mathcal F_{n,R}(T)-\mathcal F_{n,S}(T),
  \tag{3.8}
\end{equation*}
where
\begin{equation*}
  \mathcal F_{n,S}(T)
  :=
  \max_{\substack{p_i\ge0\ (i\in S)\\
                   \sum_{i\in S} p_i=1\\
                   \sum_{i\in S} i p_i=T}}
  \left[
    -\sum_{i\in S} p_i\log p_i+\sum_{i\in S} p_i h_n(i)
  \right].
  \tag{3.8a}
\end{equation*}
Fix the compact target interval
\begin{equation*}
  K_*:=
  \left[
    \frac2{\log 2}-\frac1{10},
    1+\frac2{\log 2}+\frac1{10}
  \right]
  \subset(2,5).
  \tag{3.9}
\end{equation*}
Then, uniformly for $T\in K_*$,
\begin{equation*}
  \mathcal F_{n,S}(T)
  \longrightarrow
  \mathcal F_S(T)
  :=
  \max_{\substack{p_i\ge0\ (i\in S)\\
                   \sum_{i\in S} p_i=1\\
                   \sum_{i\in S} i p_i=T}}
  \left[
    -\sum_{i\in S} p_i\log p_i
    -\frac{\log 2}{2}\sum_{i\in S} i^2p_i
  \right].
  \tag{3.9a}
\end{equation*}
If $\lambda_{n,S}(T)$ and $\lambda_S(T)$ are the finite and limiting
effective tilts, then there is a fixed $\Lambda>0$ such that
\[
  |\lambda_{n,S}(T)|+|\lambda_S(T)|\le2\Lambda
\]
for all sufficiently large $n$, and the deterministic dual error
\[
\begin{split}
  \varepsilon_n^{\mathrm{dual}}
  :=\max_{S\in\{S_+,S_4\}}
  \sup_{T\in K_*}
  \bigl(&|\mathcal F_{n,S}(T)-\mathcal F_S(T)|\\
        &+|\lambda_{n,S}(T)-\lambda_S(T)|\bigr)
\end{split}
\]
satisfies $\varepsilon_n^{\mathrm{dual}}\to0$.

For $S_4$, if $p_{n,i}(T)$ and $p_i(T)$ denote the finite and limiting
optimizers, then
\begin{equation*}
  \sup_{T\in K_*}\max_{2\le i\le5}
  |p_{n,i}(T)-p_i(T)|\longrightarrow0,
  \qquad
  \inf_{\substack{n\ \mathrm{large}\\T\in K_*}}
  \min_{2\le i\le5}p_{n,i}(T)>0.
  \tag{3.9b}
\end{equation*}
All error sequences and eventuality thresholds above are independent of
$\delta$.
\end{lemma}

\begin{proof}
Put $q:=\log 2$, $L:=\log n$, and $\ell:=\log\log n$.
Dividing (3.2) by $k=n/s$ and writing $p_i=k_i/k$ gives
\[
  \frac{L_S(n,n/s)}{n/s}
  =(s-1)L-s+1+\log s
   +\max_p\left[-\sum_i p_i\log p_i
                 -\sum_i p_i\log d_{\alpha-i}\right].
\]
There is an exact affine-plus-curved decomposition
\begin{equation*}
  -\log d_{\alpha-i}=A_n+B_ni+h_n(i),
  \tag{3.10}
\end{equation*}
where
\begin{equation*}
  A_n:=-\log d_\alpha,
  \qquad
  B_n:=q\alpha-\frac q2+\log\alpha,
  \tag{3.11}
\end{equation*}
and $h_n$ is the curved score defined in (3.1a). Direct subtraction verifies
the displayed identity also at $i=-1$.
Thus
\begin{equation*}
  h_n(i)=-\frac q2i^2+O(i^2/\alpha)
  \quad\text{for every fixed }i,
  \qquad
  h_n(i)\le-\frac q2i^2
  \quad(i\ge-1).
  \tag{3.12}
\end{equation*}
The first estimate is uniform on every fixed finite set of deficits; the
second is global on the finite support.

\displayheading{Finite duals and optimizing tilts}
For $S\in\{S_+,S_4\}$, define
\[
  Z_{n,S}(\lambda)
  :=\sum_{i\in S}
    \exp\{\lambda i+h_n(i)\},
  \qquad
  M_{n,S}(\lambda)
  :=\frac{d}{d\lambda}\log Z_{n,S}(\lambda),
\]
with the cutoff $S_+^{(n)}$ in the finite unrestricted sum. Define the
limiting functions by
\[
  Z_S(\lambda)
  :=\sum_{i\in S}
    \exp\!\left(\lambda i-\frac q2i^2\right),
  \qquad
  M_S(\lambda):=\frac{d}{d\lambda}\log Z_S(\lambda).
\]
The target center satisfies
\begin{equation*}
  \frac2q\le T_0\le1+\frac2q,
  \tag{3.13}
\end{equation*}
so it lies in the interior of $K_*$. The same is true of every target in
$K_*$ relative to both supports.

Choose $\Lambda>0$ so that
\[
  M_S(-\Lambda)<\min K_*<\max K_*<M_S(\Lambda)
  \qquad(S=S_+,S_4).
\]
Such a choice exists because the limiting mean maps are strictly increasing
and their endpoint limits bracket $K_*$. For $m=0,1,2$ and
$|\lambda|\le\Lambda$, (3.12) gives the summable majorant
\[
  |i|^m\exp\!\left(\Lambda|i|-\frac q2i^2\right).
\]
On every fixed finite set of indices, the weights converge by (3.12). The
majorant makes the remaining Gaussian tail uniformly small, and the omitted
tail beyond the finite cutoff $\alpha-1$ is bounded by the same series.
Consequently the zeroth, first, and second tilted moments converge uniformly
on $[-\Lambda,\Lambda]$. In particular,
\[
  Z_{n,S}\to Z_S,
  \qquad
  M_{n,S}\to M_S,
  \qquad
  M'_{n,S}\to M'_S
\]
uniformly there.

Now
$M'_S(\lambda)=\operatorname{Var}_{S,\lambda}(i)>0$. By continuity and
compactness,
\[
  v_*:=\frac12
  \min_{\substack{S\in\{S_+,S_4\}\\|\lambda|\le\Lambda}}
  M'_S(\lambda)>0.
\]
For all sufficiently large $n$, one has
$M'_{n,S}(\lambda)\ge v_*$ on the same interval, and the finite mean maps
also bracket $K_*$. Hence, for every $T\in K_*$, the equations
\[
  M_{n,S}(\lambda_{n,S}(T))=T,
  \qquad
  M_S(\lambda_S(T))=T
\]
have unique solutions in $[-\Lambda,\Lambda]$. The inverse mean maps are
uniformly $v_*^{-1}$-Lipschitz, and therefore
\[
  \sup_{T\in K_*}
  |\lambda_{n,S}(T)-\lambda_S(T)|
  \le
  v_*^{-1}\sup_{|\lambda|\le\Lambda}
  |M_{n,S}(\lambda)-M_S(\lambda)|
  \longrightarrow0.
\]

The entropy functional in (3.8a) is strictly concave on its feasible
simplex. Its unique optimizer is
\[
  p_{n,S,T}(i)
  =\frac{
    \exp\{\lambda_{n,S}(T)i+h_n(i)\}
  }{Z_{n,S}(\lambda_{n,S}(T))},
\]
and the dual value is exactly
\[
  \mathcal F_{n,S}(T)
  =\log Z_{n,S}(\lambda_{n,S}(T))
   -\lambda_{n,S}(T)T.
\]
The corresponding formulas hold in the limit. Uniform convergence of the
partition functions and tilts proves (3.9a) and
$\varepsilon_n^{\mathrm{dual}}\to0$. On the finite support $S_4$, the optimizer
formula also gives uniform coordinatewise convergence. The limiting
coordinates are positive and continuous in $T$ on the compact set $K_*$;
this proves (3.9b).

\displayheading{Exact cancellation between supports}
Substituting (3.10) into the normalized exponent gives the exact identity
\begin{equation*}
\begin{split}
  \Psi_{n,S}(s)
  &:=\frac{L_S(n,n/s)}{n/s}\\
  &=(s-1)L-s+1+\log s
    +A_n+B_n(\alpha-s)
    +\mathcal F_{n,S}(\alpha-s).
\end{split}
  \tag{3.14}
\end{equation*}
At a fixed feasible $k$, the values of $s$ and $T=\alpha-s$ are the same for
both supports. Every term in (3.14) except the final dual value is therefore
common, which proves the exact difference identity (3.8). No limiting
replacement is used in this cancellation.

\displayheading{The value and derivative at the phase center}
At $s=s_0$, let $T_0=\alpha-s_0$. Applying Stirling once to $A_n$ in
(3.14), and then simplifying the affine terms exactly, gives the scalar
expression
\begin{equation*}
\begin{split}
  \Psi_{n,S}(s_0)-\mathcal F_{n,S}(T_0)
  ={}&s_0\left(
    L-\log\alpha-\frac q2s_0+\frac q2
  \right)-L+T_0+1+\log s_0+\frac q2T_0^2\\
  &\hspace{42mm}
  -\frac12\log(2\pi\alpha)+O(1/L).
\end{split}
  \tag{3.15}
\end{equation*}
Equation (2.1) gives the exact identity
\[
  \frac{qs_0}{2}=L-\ell+\log q-q.
\]
Since $\alpha=s_0+T_0$ and $T_0$ is uniformly bounded,
\[
  \log\alpha
  =\log s_0+\frac{T_0}{s_0}+O(L^{-2}),
  \qquad
  \log s_0=\ell+q-\log q+O(\ell/L).
\]
Consequently
\[
\begin{split}
  s_0\left(
    L-\log\alpha-\frac q2s_0+\frac q2
  \right)
  &=-T_0+\frac q2s_0+O(\ell)\\
  &=-T_0+L-\ell+\log q-q+O(\ell).
\end{split}
\]
Substitution in (3.15) cancels the terms $L$ and $T_0$ and leaves
$O(\ell)$ uniformly in the phase. The dual term in (3.14) is $O(1)$ uniformly
on $K_*$, by the Gaussian majorant above. Therefore
\begin{equation*}
  \Psi_{n,S}(s_0)=O(\ell)
  \tag{3.16}
\end{equation*}
uniformly for both supports.

The phase formula for $\alpha$ and the expansion of $\log\alpha$ also give
\begin{equation*}
  B_n=2L-\ell+O(1)
  \tag{3.17}
\end{equation*}
uniformly in $\delta$. Strict concavity and the optimizer formula imply
\[
  \frac{d}{dT}\mathcal F_{n,S}(T)=-\lambda_{n,S}(T).
\]
Differentiating (3.14), with $T=\alpha-s$, yields
\begin{equation*}
  \Psi'_{n,S}(s)
  =L-1+s^{-1}-B_n+\lambda_{n,S}(T)
  =-L+O_A(\ell)
  \tag{3.18}
\end{equation*}
uniformly whenever $|s-s_0|\le A\ell/L$. Indeed, the image of this corridor
under $T=\alpha-s$ lies inside $K_*$ for all sufficiently large $n$, because
$K_*$ has a fixed margin around the complete range of $T_0$.

\displayheading{Existence, uniqueness, and slope of the roots}
Set
\[
  \Psi_{n,S,c}(s):=\Psi_{n,S}(s)+c.
\]
By (3.16), there is an absolute $C_0$ such that
$|\Psi_{n,S,c}(s_0)|\le C_0\ell$ for both supports and every
$0\le c\le q$. By (3.18), after one phase-independent eventuality threshold,
\[
  \Psi'_{n,S,c}(s)\le-\frac12L
\]
throughout each fixed corridor. Choose $A_*>2C_0$. Integration gives
\[
\begin{aligned}
  \Psi_{n,S,c}\!\left(s_0-A_*\frac\ell L\right)
  &\ge\left(\frac{A_*}{2}-C_0\right)\ell>0,\\
  \Psi_{n,S,c}\!\left(s_0+A_*\frac\ell L\right)
  &\le\left(C_0-\frac{A_*}{2}\right)\ell<0.
\end{aligned}
\]
The function is strictly decreasing in $s$ on this interval, so it has one and
only one zero there. Since $k=n/s>0$, this is equivalent to the unique root
$r_{S,c}$ of $L_S(n,k)+ck=0$, and proves (3.6).

Finally,
\[
  L_S(n,k)+ck=k\Psi_{n,S,c}(s),
  \qquad
  \frac{ds}{dk}=-\frac sk.
\]
Hence
\begin{equation*}
\begin{split}
  \frac{\partial}{\partial k}\{L_S(n,k)+ck\}
  &=\Psi_{n,S,c}(s)-s\Psi'_{n,S}(s)\\
  &=\frac2qL^2+O_A(L\ell),
\end{split}
  \tag{3.19}
\end{equation*}
because $s=(2/q)L+O_A(\ell)$, (3.18) holds, and
$\Psi_{n,S,c}(s)=O_A(\ell)$ throughout the corridor. This proves (3.7) and
$\varepsilon_{n,A}^{\mathrm{slope}}\to0$ with one eventuality threshold for
the complete phase.
\end{proof}
\section{\texorpdfstring{A uniform lower location for
$\chi$}{A uniform lower location for chi}}
\label{a-valid-unrestricted-lower-location-for-chi}

Let $r_+(n)=r_{S_+,0}$ be the unrestricted zero from Lemma~3.1, and put
\begin{equation*}
  k_\chi^-:=\lfloor r_+(n)\rfloor-\lceil\log n\rceil.
  \tag{4.1}
\end{equation*}
This is a deterministic integer. Since
$r_+(n)=\Theta(n/\log n)$ uniformly in the phase, one has
$1\le k_\chi^-<n$ for all sufficiently large $n$.

For an integer profile
$\mathbf k=(k_i)_{-1\le i\le\alpha-1}$, let
$X_{\mathbf k}$ be the number of unordered proper colorings having exactly
$k_i$ classes of size $\alpha-i$. Direct enumeration gives
\begin{equation*}
  \Exp{X_{\mathbf k}}
  =\frac{n!}{\prod_i((\alpha-i)!)^{k_i}k_i!}
   2^{-\sum_i k_i\binom{\alpha-i}{2}}.
  \tag{4.2}
\end{equation*}
For an integer $k$, let $\mathcal K_{n,k,\alpha+1}$ be the set of these
profiles satisfying
\[
  \sum_i k_i=k,
  \qquad
  \sum_i(\alpha-i)k_i=n,
\]
and define
\[
  E_{n,k,\alpha+1}
  :=\sum_{\mathbf k\in\mathcal K_{n,k,\alpha+1}}
    \Exp{X_{\mathbf k}}.
\]
Thus $E_{n,k,\alpha+1}$ is the expected number of unordered proper colorings
with exactly $k$ nonempty parts, each of size at most $\alpha+1$.

\subsection*{The profile sum}

There are $\alpha+1=O(\log n)$ profile coordinates, and each coordinate lies
between zero and $n$. Hence
\begin{equation*}
  |\mathcal K_{n,k,\alpha+1}|
  \le(n+1)^{\alpha+1}
  =\exp\{O((\log n)^2)\}
  \tag{4.3a}
\end{equation*}
uniformly in $k$ and in the phase. Applying the uniform factorial estimate
(1.2) to (4.2), with the convention $0\log0=0$, gives
\[
  \log\Exp{X_{\mathbf k}}
  \le L_{\mathbf k}+O((\log n)^2).
\]
The error is uniform even when some $k_i=0$: only $O(\log n)$ factorials
occur, and every one contributes $O(\log n)$. Maximizing over the feasible
profile and then summing (4.3a) yields
\begin{equation*}
  \log E_{n,k,\alpha+1}
  \le L_+(n,k)+O((\log n)^2).
  \tag{4.3}
\end{equation*}

\subsection*{Displacement below the root}

Set
\[
  \Delta_n:=r_+(n)-k_\chi^-.
\]
The floor and ceiling in (4.1) give the deterministic bounds
\begin{equation*}
  \log n\le\Delta_n<\log n+2.
  \tag{4.3b}
\end{equation*}
We next verify that the entire interval
$[k_\chi^-,r_+(n)]$ lies in the derivative corridor of Lemma~3.1. Write
$s(k)=n/k$. Uniformly on this interval,
$k=\Theta(n/\log n)$ and $s(k)=\Theta(\log n)$. Therefore
\[
  |s(k_\chi^-)-s(r_+)|
  \le
  \sup_{k\in[k_\chi^-,r_+]}
  \frac{n}{k^2}\,\Delta_n
  =O\!\left(\frac{(\log n)^3}{n}\right)
  =o\!\left(\frac{\log\log n}{\log n}\right).
\]
Thus, after one phase-independent threshold, every point of the interval lies
inside the common corridor used in (3.7).

Equation (3.7), with $S=S_+$ and $c=0$, consequently gives an absolute
$c_*>0$ such that
\begin{equation*}
  \frac{d}{dk}L_+(n,k)\ge c_*(\log n)^2
  \qquad
  (k_\chi^-\le k\le r_+)
  \tag{4.3c}
\end{equation*}
for all sufficiently large $n$, uniformly in the phase. Since
$L_+(n,r_+)=0$, the mean-value theorem and (4.3b) imply
\begin{equation*}
  L_+(n,k_\chi^-)
  \le-c_*(\log n)^2\Delta_n
  \le-c_*(\log n)^3.
  \tag{4.4}
\end{equation*}
Combining (4.3) and (4.4), and decreasing the constant once, gives a
deterministic sequence
\[
  \varepsilon_n^{\mathrm{prof}}
  :=
  \exp\{-c_\chi(\log n)^3\}
  \longrightarrow0
\]
such that
\begin{equation*}
  E_{n,k_\chi^-,\alpha+1}
  \le\varepsilon_n^{\mathrm{prof}}.
  \tag{4.4a}
\end{equation*}

\subsection*{Removing the size cap}

Let
\[
  \mathcal A_n:=\{\alpha(G_n)\le\alpha+1\}.
\]
Equation (2.9) gives a deterministic phase-uniform sequence
$\varepsilon_n^{\mathrm{cap}}\to0$ such that
\[
  \Prob{\mathcal A_n^c}
  \le\varepsilon_n^{\mathrm{cap}}.
\]
On $\mathcal A_n$, every class of every proper coloring has size at most
$\alpha+1$.

Suppose that $\chi(G_n)\le k_\chi^-$ and choose a proper coloring with
$h\le k_\chi^-$ nonempty classes. If $h<k_\chi^-$, then some class contains at
least two vertices: otherwise $h=n$, contradicting $h<k_\chi^-<n$. Splitting
such a class into two nonempty subsets preserves independence. Repeating this
operation produces a proper coloring with exactly $k_\chi^-$ nonempty classes.
The size cap is preserved under splitting. Therefore
\[
  \{\chi(G_n)\le k_\chi^-\}\cap\mathcal A_n
  \subseteq
  \{\text{an exactly $k_\chi^-$, $(\alpha+1)$-bounded coloring exists}\}.
\]
Markov's inequality and (4.4a) now give
\begin{equation*}
\begin{aligned}
  \Prob{\chi(G_n)\le k_\chi^-}
  &\le
  \Prob{\mathcal A_n^c}
  +E_{n,k_\chi^-,\alpha+1}\\
  &\le
  \varepsilon_n^{\mathrm{cap}}
  +\varepsilon_n^{\mathrm{prof}}
  \longrightarrow0.
\end{aligned}
  \tag{4.5}
\end{equation*}
Equivalently,
\begin{equation*}
  \Prob{\chi(G_n)>k_\chi^-}\longrightarrow1.
  \tag{4.6}
\end{equation*}
The direction is strict, as required by the final event intersection.

Finally, (4.3b) also records the location precision:
\begin{equation*}
  |k_\chi^--r_+(n)|
  <\log n+2
  =o\!\left(\frac{n}{(\log n)^3}\right).
  \tag{4.7}
\end{equation*}
Thus both the probability statement and the location error hold with one set
of phase-independent thresholds along the full sequence of integers.

\section{The four-size signed first-moment advantage}\label{the-four-size-signed-first-moment-advantage}

Recall that a signed cocoloring witness is a profile partition with one
independent-or-complete declaration on each class. It is realized when each
class induces the declared graph. These declarations are auxiliary counting
data, not a new graph invariant.

The four-size comparison has two steps. Restricting the deficits to
\(S_4=\{2,3,4,5\}\) must cost strictly less than \(\log 2\) per part. The
\(2^k\) sign choices then convert this strict entropy margin into a
macroscopic root separation. Lemma~5.1 proves the margin, and the slope
estimate in Lemma~3.1 converts it into displacement.

For \(S\in\{S_+,S_4\}\), introduce the tilted weights, partition
function, and mean map

\[
 w_i(\lambda):=\exp\!\left(\lambda i-\frac{\log 2}{2}i^2\right),
 \qquad
 Z_S(\lambda):=\sum_{i\in S}w_i(\lambda),
\]

\[
 M_S(\lambda):=
 \frac{\sum_{i\in S}i\,w_i(\lambda)}{Z_S(\lambda)}
 =\frac{d}{d\lambda}\log Z_S(\lambda).
\]

The Gaussian factor makes both sums finite.  Moreover,
\(M_S'(\lambda)=\operatorname{Var}_{S,\lambda}(i)>0\), so for every
target mean \(T\) in the interior of the convex hull of \(S\) there is a
unique tilt \(\lambda_S(T)\) with \(M_S(\lambda_S(T))=T\).  The optimizer
in (3.9a) is therefore

\[
 p_i=\frac{w_i(\lambda_S(T))}{Z_S(\lambda_S(T))}
 =\frac{e^{\lambda_S(T)i-{\log 2}i^2/2}}
        {\sum_{j\in S}e^{\lambda_S(T)j-{\log 2}j^2/2}}.       \tag{5.1}
\]

The same calculation gives the dual representation

\[
 \mathcal F_S(T)
 =\inf_{\lambda\in\mathbb R}
   \{\log Z_S(\lambda)-\lambda T\}
 =\log Z_S(\lambda_S(T))-\lambda_S(T)T.
\]

Define

\[
 D_4(\delta)=\mathcal F_{S_+}(T_0)
              -\mathcal F_{S_4}(T_0),\qquad
 \gamma_4=\log\frac{200}{153}.                            \tag{5.2}
\]
This certificate is sufficient both for the root separation and for the
displayed numerical constant in the final theorem.

\begin{lemma}[Uniform entropy certificate]
\label{lem:entropy-gap}

For every \(0\le \delta\le 1\),

\[
 0\le D_4(\delta)<\log\frac{153}{100},\qquad
 {\log 2}-D_4(\delta)>\gamma_4.                                 \tag{5.3}
\]

\end{lemma}

\begin{proof}
Fix \(\delta\in[0,1]\), put
\(T_0=1+2/{\log 2}-\delta\), and abbreviate

\[
 \lambda_4:=\lambda_{S_4}(T_0).
\]

We begin by bracketing this tilt. At \(\lambda=2{\log 2}\), set
\(j=i-2\), a bijection from \(S_4\) onto \(\{0,1,2,3\}\).
Substituting \(i=j+2\) gives

\[
\begin{aligned}
 w_{j+2}(2{\log 2})
 &=\exp\!\left(2{\log 2}(j+2)
       -\frac{{\log 2}}2(j+2)^2\right)\\
 &=\exp\!\left(2{\log 2}-\frac{{\log 2}}2j^2\right)
 =4\,2^{-j^2/2}.
\end{aligned}
\]

Consequently, reindexing both sums in the definition of the mean---and
only reindexing them, without discarding any term---gives

\[
\begin{aligned}
 Z_{S_4}(2{\log 2})
 &=\sum_{i=2}^{5}w_i(2{\log 2})
   =4\sum_{j=0}^{3}2^{-j^2/2},\\
 \sum_{i=2}^{5}i\,w_i(2{\log 2})
 &=4\sum_{j=0}^{3}(j+2)2^{-j^2/2}\\
 &=8\sum_{j=0}^{3}2^{-j^2/2}
   +4\sum_{j=0}^{3}j2^{-j^2/2}.
\end{aligned}
\]

Dividing the second equality by the first explains the additive
constant \(2\) in the following formula:

\[
 \begin{aligned}
 M_{S_4}(2{\log 2})
 &=2+\frac{\sum_{j=0}^{3}j2^{-j^2/2}}
          {\sum_{j=0}^{3}2^{-j^2/2}}\\
 &<2+\frac{\sum_{j\ge0}j2^{-j^2/2}}
          {\sum_{j\ge0}2^{-j^2/2}}
 <2+\frac45<\frac2{\log 2}\le T_0.
 \end{aligned}
 \tag{5.4}
\]

We verify the first strict inequality in (5.4) directly. Put
\[
 A=\sum_{j=0}^{3}2^{-j^2/2},\quad
 B=\sum_{j=0}^{3}j2^{-j^2/2},\quad
 C=\sum_{j=4}^{\infty}2^{-j^2/2},\quad
 D=\sum_{j=4}^{\infty}j2^{-j^2/2}.
\]
Here \(B/A\le3\), while \(D/C\ge4\).  Since \(A,C>0\),
\[
 \frac{B+D}{A+C}-\frac BA
 =\frac{AD-BC}{A(A+C)}>0.
\]
Thus adjoining the terms whose new indices satisfy \(j\ge4\) strictly
increases the weighted mean. For the next inequality, put
\(a_j=j2^{-j^2/2}\). For \(j\ge4\),
\[
 \frac{a_{j+1}}{a_j}
 =\frac{j+1}{j}\,2^{-(2j+1)/2}\le\frac1{16},
\]
and \(a_4=1/64\). Hence
\[
 \sum_{j\ge4}j2^{-j^2/2}
 \le\frac{1/64}{1-1/16}=\frac1{60}.
\]
Writing \(b=2^{-1/2}\), retaining \(j=0,1,2\) in the denominator and
\(j=0,1,2,3\) in the numerator gives
\[
 \frac{\sum_{j\ge0}j2^{-j^2/2}}
      {\sum_{j\ge0}2^{-j^2/2}}
 <\frac{19b/16+1/2+1/60}{5/4+b}<\frac45,
\]
where the last inequality follows from \(b<71/100\). Finally,
\(2/3<\log 2<5/7\) gives \(2+4/5<2/\log 2\).

At \(\lambda=9{\log 2}/2\), put \(t=2^{-1/8}\). The identity
\[
 \frac{9\log 2}{2}i-\frac{\log 2}{2}i^2
 =\frac{\log 2}{8}\{81-(2i-9)^2\}
\]
gives
\[
 w_i(9{\log 2}/2)=2^{81/8}t^{(2i-9)^2}.
\]
Thus the four weights, in the order
\(i=2,3,4,5\), are proportional to

\[
 t^{25},\quad t^9,\quad t,\quad t.
\]

The numerator of \(M_{S_4}(9{\log 2}/2)-4\) is therefore

\[
 t-t^9-2t^{25}=t(1-t^8-2t^{24})=t/4>0,                 \tag{5.5}
\]

and the denominator is positive.  Hence this mean is greater than four,
whereas \(T_0\le1+2/{\log 2}<4\).  Since \(M_{S_4}\) is strictly
increasing, the two mean comparisons give

\[
 2{\log 2}<\lambda_4<9{\log 2}/2.                                     \tag{5.6}
\]

For a tilt \(\lambda\), define the omitted low- and high-deficit ratios

\[
 L(\lambda):=
 \frac{\sum_{i=-1}^{1}w_i(\lambda)}{Z_{S_4}(\lambda)},
 \qquad
 H(\lambda):=
 \frac{\sum_{i\ge6}w_i(\lambda)}{Z_{S_4}(\lambda)}.
\]

Let \(M_A(\lambda)\) denote the mean of \(i\) under the weights
\(w_i(\lambda)\) restricted to an index set \(A\). Differentiation gives
\[
 \frac{d}{d\lambda}\log L(\lambda)
 =M_{\{-1,0,1\}}(\lambda)-M_{S_4}(\lambda)
 \le 1-2=-1,
\]
and
\[
 \frac{d}{d\lambda}\log H(\lambda)
 =M_{\{6,7,\ldots\}}(\lambda)-M_{S_4}(\lambda)
 \ge 6-5=1.
\]
Thus \(L\) is decreasing and \(H\) is increasing on \(\mathbb R\).
Direct summation and Gaussian-tail bounds give

\[
 L(2{\log 2})<\frac{51}{100},\quad H(3{\log 2})<\frac1{50},\quad
 L(3{\log 2})<\frac3{25},\quad H(9{\log 2}/2)<\frac14.                 \tag{5.7}
\]

Here are explicit checks. At \(2{\log 2}\), with \(b=2^{-1/2}\),

\[
 L(2{\log 2})=\frac{b+1/4+b/16}{1+b+1/4+b/16}<51/100.
\]

At \(\lambda=3\log 2\),
\[
 w_i(3\log 2)
 =\exp\!\left((\log 2)\left(3i-\frac{i^2}{2}\right)\right)
 =2^{9/2}2^{-(i-3)^2/2}.
\]
Set \(b=2^{-1/2}\). After canceling the common factor \(2^{9/2}\),
the weight at index \(i\) is \(b^{(i-3)^2}\). Thus the four
\(S_4\)-weights, for \(i=2,3,4,5\), are
\[
 b,\quad 1,\quad b,\quad b^4,
\]
and
\[
 2^{-9/2}Z_{S_4}(3\log 2)=1+2b+b^4=\frac54+2b.
\]
For the omitted low indices \(i=-1,0,1\), the weights are
\[
 b^{16},\quad b^9,\quad b^4
 =\frac1{256},\quad\frac b{16},\quad\frac14.
\]
For the omitted high indices \(i\ge6\), they are
\[
 b^9+b^{16}+\sum_{q\ge5}b^{q^2}.
\]
The first term in the remaining sum is \(b^{25}=b/4096\), and the
ratio between successive terms is at most \(b^{11}<1/32\). Therefore
\[
 \sum_{q\ge5}b^{q^2}
 \le\frac{b^{25}}{1-1/32}<\frac1{3968}.
\]
It follows that
\[
 L(3\log 2)
 =\frac{1/256+b/16+1/4}{5/4+2b}<\frac3{25},
\]
and
\[
 H(3\log 2)
 \le\frac{b/16+1/256+1/3968}{5/4+2b}<\frac1{50},
\]
where the final inequalities follow from \(7/10<b<71/100\).

Finally, at \(9\log 2/2\) the normalized exponent is
\((2i-9)^2\). Dividing all weights by the common \(S_4\) weight \(t\),
the high indices \(i=6,7,8\) contribute \(t^8,t^{24},t^{48}\).
From \(i=8\) onward the ratio between successive terms is at most
\(1/16\), so the tail beginning with \(t^{48}\) is less than \(1/60\).
The normalized \(S_4\) denominator is \(2+t^8+t^{24}\). Hence
\(H(9\log 2/2)<1/4\) follows from

\[
 3(t^8+t^{24})+4/60<2,                                  \tag{5.8}
\]

because \(t^8=1/2\) and \(t^{24}=1/8\).

If \(\lambda_4\le 3{\log 2}\), monotonicity and (5.7) give
\(L(\lambda_4)+H(\lambda_4)<53/100\); if
\(\lambda_4\ge 3{\log 2}\), they give
\(L(\lambda_4)+H(\lambda_4)<37/100\).
By the dual representation above,

\[
 \mathcal F_{S_4}(T_0)
 =\log Z_{S_4}(\lambda_4)-\lambda_4T_0,
\]

whereas evaluating the \(S_+\) dual function at the same parameter gives

\[
 \mathcal F_{S_+}(T_0)
 \le \log Z_{S_+}(\lambda_4)-\lambda_4T_0.
\]

Subtracting these two relations yields

\[
 \begin{aligned}
 D_4(\delta)
 &\le \log\frac{Z_{S_+}(\lambda_4)}{Z_{S_4}(\lambda_4)}\\
 &=\log\!\bigl(1+L(\lambda_4)+H(\lambda_4)\bigr)
 <\log(153/100).
 \end{aligned}
\]

Since \(S_4\) is a subset of \(S_+\), the loss is
nonnegative. Subtracting from \(\log 2\) completes the proof.
\end{proof}

The following \(2^k\) gain is the one-partition identity isolated by
\citet[Proposition~6]{heckel-2025-difference}.  We repeat its short
argument because it is an input to the new root calculation.  For a fixed
partition into \(k\) classes, assigning each class the declaration
``independent'' or ``complete'' multiplies its first moment by \(2^k\):
every declaration prescribes all internal edge bits, and all \(2^k\)
declarations have probability \(2^{-\sum \binom{u}{2}}\).  The declarations
are disjoint here because all four allowed class sizes tend to infinity
(in particular, there are no singleton classes).
Let
\(r_4^{\mathrm{co}}=r_{S_4,{\log 2}}\)
be the unique corridor zero of

\[
 L_{S_4}(n,k)+{\log 2}k.                                       \tag{5.9}
\]



Put
\[
  \Phi_n(k):=L_{S_4}(n,k)+(\log 2)k.
\]
Then $\Phi_n(r_4^{\mathrm{co}})=0$ by definition. To evaluate the same
function at the unrestricted root, let
\[
  T_+(n):=\alpha-\frac{n}{r_+(n)}
\]
and introduce the finite-$n$ support loss
\[
  D_{4,n}(T):=
  \mathcal F_{n,S_+}(T)-\mathcal F_{n,S_4}(T).
\]
Since $L_{S_+}(n,r_+)=0$, the exact support-comparison identity (3.8) gives
\begin{equation*}
\begin{aligned}
  \Phi_n(r_+)
  &=L_{S_4}(n,r_+)-L_{S_+}(n,r_+)+(\log 2)r_+\\
  &=r_+\{\log 2-D_{4,n}(T_+(n))\}.
\end{aligned}
\tag{5.9a}
\end{equation*}
This equality is finite and contains no limiting substitution.

We next compare the finite support loss at $T_+(n)$ with the limiting support
loss at the phase-center target $T_0$. Define the deterministic target
displacement
\[
  \varepsilon_n^{\mathrm{target}}
  :=|T_+(n)-T_0|.
\]
Equation (3.6) gives the equivalent phase-uniform statements
\begin{equation*}
  T_+(n)=T_0+O\!\left(\frac{\log\log n}{\log n}\right),
  \qquad
  \varepsilon_n^{\mathrm{target}}
  =O\!\left(\frac{\log\log n}{\log n}\right)
  \longrightarrow0.
  \tag{5.9b}
\end{equation*}
The common corridor lies inside the fixed compact interval $K_*$ for all
sufficiently large $n$.

The limiting dual value is differentiable, with
\[
  \mathcal F_S'(T)=-\lambda_S(T),
\]
because the envelope theorem cancels the derivative of the optimizing tilt.
Section~3 bounds the limiting tilts by one constant on $K_*$, so the two
limiting dual values have a common Lipschitz constant $C_{\mathrm{Lip}}$.
The deterministic finite-dual error from Section~3 satisfies
\[
  \sup_{T\in K_*}
  |\mathcal F_{n,S}(T)-\mathcal F_S(T)|
  \le\varepsilon_n^{\mathrm{dual}}
  \qquad(S=S_+,S_4).
\]
Consequently the single theorem-facing root error
\begin{equation*}
  \omega_n^{\mathrm{root}}
  :=2\varepsilon_n^{\mathrm{dual}}
    +2C_{\mathrm{Lip}}\varepsilon_n^{\mathrm{target}}
  \longrightarrow0
  \tag{5.9c}
\end{equation*}
is deterministic and phase-uniform. Indeed,
\[
\begin{aligned}
  |D_{4,n}(T_+(n))-D_4(\delta)|
  \le{}&
  |\mathcal F_{n,S_+}(T_+)-\mathcal F_{S_+}(T_+)|\\
  &+|\mathcal F_{n,S_4}(T_+)-\mathcal F_{S_4}(T_+)|\\
  &+|\mathcal F_{S_+}(T_+)-\mathcal F_{S_+}(T_0)|\\
  &+|\mathcal F_{S_4}(T_+)-\mathcal F_{S_4}(T_0)|\\
  \le{}&\omega_n^{\mathrm{root}}.
\end{aligned}
\]
Thus $\omega_n^{\mathrm{root}}$ is one deterministic error sequence valid
across the complete phase, including integer sequences approaching either
endpoint.

Substitution into (5.9a) yields the theorem-facing estimate
\begin{equation*}
  L_{S_4}(n,r_+)+(\log 2)r_+
  =r_+\{\log 2-D_4(\delta)+O(\omega_n^{\mathrm{root}})\}
  =r_+\{\log 2-D_4(\delta)+o(1)\}.
  \tag{5.10}
\end{equation*}
Thus the $o(1)$ in (5.10) is not an independent asymptotic convention: it is
controlled by the explicit sum of the finite-dual and target-displacement
errors.

Finally, Lemma~5.1 gives
$\log 2-D_4(\delta)\ge\gamma_4>0$. Hence (5.10) implies
$\Phi_n(r_+)>0$ for all sufficiently large $n$, uniformly in the phase. The
normalized slope error in (3.7) tends to zero, so the derivative of $\Phi_n$
is positive throughout the common root corridor after one phase-independent
threshold. Since $\Phi_n(r_4^{\mathrm{co}})=0$, this proves
\[
  r_4^{\mathrm{co}}<r_+.
\]
The interval used in the subsequent mean-value argument is therefore
nonempty, with its orientation fixed.
Both roots lie in the corridor (3.6), so the whole interval between them
is inside a fixed corridor.  The mean-value theorem therefore gives a
point \(\xi_n\in(r_4^{\mathrm{co}},r_+)\) such that
\[
 \Phi_n(r_+)-\Phi_n(r_4^{\mathrm{co}})
 =\Phi_n'(\xi_n)(r_+-r_4^{\mathrm{co}}).
\]
Now \(r_+=({\log 2}/2+o(1))n/{\log n}\), and (3.7), uniformly at
\(\xi_n\), gives
\[
 \Phi_n'(\xi_n)=\frac2{\log 2}(\log n)^2
                 +O((\log n)(\log\log n)).
\]
Substituting this and (5.10) into the preceding identity yields

\[
 r_+-r_4^{\mathrm{co}}
 =\left(\frac{(\log 2)^2}{4}\{{\log 2}-D_4(\delta)\}+o(1)\right)
   \frac n{(\log n)^3}.                                        \tag{5.11}
\]

In particular, for all sufficiently large \(n\),

\[
 r_+-r_4^{\mathrm{co}}\ge\frac{(\log 2)^2\gamma_4}{8}\frac n{(\log n)^3}.     \tag{5.12}
\]

Choose the midpoint integer

\[
 k_{\mathrm{co}}=\left\lceil\frac{r_4^{\mathrm{co}}+r_+}{2}\right\rceil. \tag{5.13}
\]

We call the resulting profile the \emph{exact midpoint profile}: its total
number of parts is the rounded midpoint between the signed four-size root
\(r_4^{\mathrm{co}}\) and the unrestricted chromatic root \(r_+\).  The
type counts are chosen next so that both finite-\(n\) conservation
constraints hold exactly.

At this exact \(k\), let \(p_i^{(n)}\) be the exact
finite-\(n\) maximizer of \(L_{S_4}\).  Let
\(\lambda_{n,\mathrm{def}}\) denote the Lagrange multiplier for the
deficit-mean constraint, with the sign convention that its contribution
is \(e^{\lambda_{n,\mathrm{def}}i}\).  The Lagrange equations and the
decomposition (3.10) then give

\[
 p_i^{(n)}\ \propto\ d_{\alpha-i}^{-1}e^{\lambda_{n,\mathrm{def}}i}
 \ \propto\ e^{\widehat\lambda_ni+h_n(i)},\qquad
 \widehat\lambda_n:=B_n+\lambda_{n,\mathrm{def}}.          \tag{5.14}
\]

Since
$e^{B_ni}e^{\lambda_{n,\mathrm{def}}i}
=e^{(B_n+\lambda_{n,\mathrm{def}})i}=e^{\widehat\lambda_ni}$,
the quantity \(\lambda_{n,\mathrm{def}}\) is the multiplier in the exact
finite-$n$ coordinates, whereas \(\widehat\lambda_n\) is the tilt in the
limiting Gaussian coordinates. (If a size-coordinate
multiplier \(\tau_n\) is used instead, then
\(\lambda_{n,\mathrm{def}}=-\tau_n\).)
The target deficit mean is \(\alpha-n/k_{\mathrm{co}}\). Lemma 3.1 makes these proportions
converge uniformly to (5.1), with
\(\widehat\lambda_n\to\lambda_{S_4}(T_0)\). Equation (5.6) also shows
directly that the four limiting weights differ by at most a fixed
factor. Thus

\[
 \min_{2\le i\le5}p_i^{(n)}\ge c_p>0                    \tag{5.15}
\]

uniformly in the phase.

Round \(k_{\mathrm{co}}p_i^{(n)}\) to preliminary integers \(\widetilde k_i\).
Define the signed constraint errors

\[
 e_0=\sum_i\widetilde k_i-k_{\mathrm{co}},\qquad
 e_1=\sum_i i\widetilde k_i-(\alpha k_{\mathrm{co}}-n),
\]

so \(e_0,e_1=O(1)\), and add

\[
 \Delta k_2=e_1-3e_0,\qquad
 \Delta k_3=2e_0-e_1.                                   \tag{5.16}
\]

Set $\Delta k_4=\Delta k_5=0$, $k_i:=\widetilde k_i+\Delta k_i$ for
$2\le i\le5$, and $\mathbf k:=(k_2,k_3,k_4,k_5)$.

Indeed, \(\Delta k_2+\Delta k_3=-e_0\) and
\(2\Delta k_2+3\Delta k_3=-e_1\). This enforces both constraints
exactly, changes only \(O(1)\) counts, and preserves positivity by
(5.15).  Put
\(\mathbf k^*=(k_{\mathrm{co}}p_i^{(n)})_{i=2}^5\) and let
\(\Delta^{\mathrm{tot}}=\mathbf k-\mathbf k^*\) be the total rounding
and correction displacement.  It has bounded norm and satisfies both
homogeneous constraint equations, because \(\mathbf k^*\) and
\(\mathbf k\) obey the same two constraints.  Thus it is tangent to the
feasible affine plane.  At the exact optimizer the linear term vanishes, while
the Hessian of the entropy term is
\(-\operatorname{diag}(1/k_i)\).  Hence

\[
 L_{\mathbf k^*+\Delta^{\mathrm{tot}}}
 -L_{\mathbf k^*}
 =O\!\left(\sum_i\frac{(\Delta_i^{\mathrm{tot}})^2}{k_i^*}\right)
 =O(1/k_{\mathrm{co}}).
\]

The resulting exact integer vector satisfies

\[
 u_i=\alpha-i,\qquad k_i=\Theta(n/{\log n})\quad(2\le i\le5).   \tag{5.17}
\]

Applying (1.2) to the factorials in the exact signed first moment of this
vector gives

\[
 \Exp{Z_{\mathbf k}^{\mathrm{sgn}}}
 =2^{k_{\mathrm{co}}}\frac{n!}{\prod_i(u_i!)^{k_i}k_i!}
   2^{-\sum_i k_i\binom{u_i}{2}},                        \tag{5.18}
\]

and

\[
 \log \Exp{Z_{\mathbf k}^{\mathrm{sgn}}}
 =L_{S_4}(n,k_{\mathrm{co}})+{\log 2}k_{\mathrm{co}}+O({\log n})
 \ge c_Zk_{\mathrm{co}}                                          \tag{5.19}
\]

for a phase-independent \(c_Z>0\). Indeed, the midpoint is at distance
at least \((r_+-r_4^{\mathrm{co}})/2-O(1)\) to the right of
\(r_4^{\mathrm{co}}\).  Integrating the positive derivative (3.7) over
that interval and using (5.12) gives the scale calculation
\[
 \frac{n}{(\log n)^3}\cdot(\log n)^2
 \asymp\frac n{\log n}\asymp k_{\mathrm{co}}.
\]
Consequently,
\[
 L_{S_4}(n,k_{\mathrm{co}})+{\log 2}k_{\mathrm{co}}
 \ge c_Zk_{\mathrm{co}},
\]
after decreasing \(c_Z\) if necessary. Every
\(u_i\) tends to infinity, so the independent and complete
declarations of one part are disjoint. Thus
\(Z_{\mathbf k}^{\mathrm{sgn}}>0\)
implies an actual cocoloring with \(k_{\mathrm{co}}\) parts.

Finally, (4.1), (5.12), and (5.13) give

\[
 k_\chi^- -k_{\mathrm{co}}
 \ge\left(\frac{(\log 2)^2\gamma_4}{16}-o(1)\right)
       \frac n{(\log n)^3}.                                    \tag{5.20}
\]

\section{Exact signed second-moment representation}\label{exact-signed-second-moment-representation}

For two partitions sharing exactly \(\ell\) whole classes,
\citet[Proposition~6]{heckel-2025-difference} bounded the number of joint
sign declarations by \(2^{2k-\ell}\), thereby transferring selected
ordinary-coloring overlap bounds to cocolorings.  Lemma 6.1 below does
not import that inequality.  It keeps the entire overlap matrix and counts
the compatible signs exactly: compatibility is controlled by the
components of the graph of cells of size at least two.  Whole common
classes appear as one special case of this component formula.

Temporarily label all slots within each of the four types. This
multiplies the unordered witness by the deterministic factor
\(\prod_i k_i!\) and hence does not change its normalized second
moment. Choose two independent uniform ordered partitions of the
profile. If their slots are indexed by \(a\) and \(b\), put

\[
 r_{ab}=|V_a\cap W_b|.                                   \tag{6.1}
\]

The overlap matrix has the exact law

\[
 p(r)=\frac{\prod_a s_a!\prod_b t_b!}
            {n!\prod_{a,b}r_{ab}!},                      \tag{6.2}
\]

where here both degree lists are the multiset given by (5.17).
Equivalently, place \(s_a\) stubs at each row slot and
\(t_b\) stubs at each column slot and take a uniform perfect
matching of the two \(n\)-stub sets.

Let \(H(r)\) be the simple bipartite graph whose edges are cells
with \(r_{ab}\ge 2\), omitting isolated slot vertices.
Write

\[
 W=\sum_{a,b}\binom{r_{ab}}2,\qquad
 \beta(H)=|E(H)|-|V(H)|+c(H).                             \tag{6.3}
\]

\begin{lemma}[Exact sign sum]
\label{lemma-6.1-exact-sign-sum}

The normalized local factor of the pair of partitions is

\[
 A_\zeta(r)=2^{W+c(H)-|V(H)|}
 =\left(\prod_{a,b}g(r_{ab})\right)2^{\beta(H)},          \tag{6.4}
\]

where

\[
 g(0)=g(1)=g(2)=1,\qquad
 g(x)=2^{\binom x2-1}\quad(x\ge3).                       \tag{6.5}
\]

Consequently

\[
 \frac{\Exp{(Z_{\mathbf k}^{\mathrm{sgn}})^2}}
      {\Exp{Z_{\mathbf k}^{\mathrm{sgn}}}^2}
 =\sum_{r\in\mathcal R(\mathbf s,\mathbf t)}
      p(r)A_\zeta(r).                                    \tag{6.6}
\]

where \(\mathcal R(\mathbf s,\mathbf t)\) is the finite set of
nonnegative integer matrices with the prescribed row margins \(\mathbf s\)
and column margins \(\mathbf t\).

\end{lemma}

\begin{proof}
Let
\(B=\sum_a \binom{s_a}{2}\),
which is the number of internal edge bits prescribed by one partition. A
row sign and a column sign are compatible on a cell of size at least two
exactly when they agree. Therefore signs must be constant on every
component of \(H\). There are

\[
 2^{2k_{\mathrm{co}}-|V(H)|+c(H)}
\]

compatible pairs of sign assignments, including the free signs on slots
isolated from \(H\). For each compatible pair, \(W\) edge bits
are prescribed twice, so the probability of all constraints is
\(2^{-(2B-W)}\). Division by the square of the one-partition
signed probability
\((2^{k_{\mathrm{co}}}2^{-B})^2\)
gives

\[
 \frac{2^{2k_{\mathrm{co}}-|V(H)|+c(H)}2^{-(2B-W)}}
      {(2^{k_{\mathrm{co}}}2^{-B})^2}
 =2^{W+c(H)-|V(H)|},
\]

which is the first expression in (6.4).

Formula (6.4) separates two effects that will be estimated by different
methods. The product of the \(g(r_{ab})\) is purely local, with
$g(0)=g(1)=g(2)=1$; the factor \(2^{\beta(H)}\) is topological and counts the
independent cycle choices of the support graph.

For every support edge, split \(\binom{r}{2}\) as
\((\binom{r}{2}-1)+1\). Since
\(|E|-|V|+c=\beta\),

\[
 \prod_{a,b}g(r_{ab})=2^{W-|E(H)|},
\]

and multiplying by \(2^{\beta(H)}\) gives
the second expression. Averaging over (6.2) proves (6.6).
\end{proof}

The identity

\[
 2^{\beta(H)}
 =\#\{F\subseteq E(H):\deg_F(v)\text{ is even for every }v\} \tag{6.7}
\]

will be used repeatedly. It is the elementary fact that the even
subgraphs form the binary cycle space of dimension \(\beta(H)\).

We also need one exact configuration-model estimate.

\begin{lemma}[Joint prescribed-cell bound]
\label{lemma-6.2-joint-prescribed-cell-bound}

In a bipartite configuration model of total degree \(m_0\), with
row degrees \(d_a\) and column degrees
\(d'_b\), prescribe distinct cells \(D\) and
demands \(x_{ab}\ge 1\). Put

\[
 x=\sum_Dx_{ab},\quad D_a=\sum_bx_{ab},\quad
 D'_b=\sum_ax_{ab}.
\]

In the feasible case

\[
 x\le m_0,\qquad D_a\le d_a\ \ (\forall a),\qquad
 D'_b\le d'_b\ \ (\forall b),
\]

we have

\[
 \Prob{r_{ab}\ge x_{ab}\ \text{for every }ab\in D}
 \le\frac{\prod_a(d_a)_{D_a}\prod_b(d'_b)_{D'_b}}
          {(m_0)_x\prod_{ab\in D}x_{ab}!}.               \tag{6.8}
\]

If any feasibility condition fails, the prescribed event is empty and is
handled separately. In particular, in all cases with \(m_0>0\), with
\(\mathrm{e}:=\exp(1)\) and
\(\theta_{ab}=\mathrm{e}\, d_a d'_b/m_0\),
its probability is at most

\[
 \prod_{ab\in D}\frac{\theta_{ab}^{x_{ab}}}{x_{ab}!}.  \tag{6.9}
\]

\end{lemma}

\begin{proof}
If a feasibility condition fails, no matching realizes all demands. In the
feasible case, assigning disjoint row stubs to the demanded cells gives
\[
  \prod_a\frac{(d_a)_{D_a}}{\prod_b x_{ab}!}
\]
choices. Independently selecting the column stubs gives
\[
  \prod_b\frac{(d'_b)_{D'_b}}{\prod_a x_{ab}!}
\]
choices. For each cell \(ab\), the selected row and column stubs can be paired
in \(x_{ab}!\) ways. Hence the product over all cells cancels one of the two
copies of \(\prod_{ab\in D}x_{ab}!\), and the total number of demand witnesses
is

\[
 \frac{\prod_a(d_a)_{D_a}\prod_b(d'_b)_{D'_b}}
      {\prod_{ab\in D}x_{ab}!}.
\]

Any fixed witness occurs with probability \(1/(m_0)_x\), because after
its \(x\) prescribed pairs are exposed the remaining matching is uniform
on the unused stubs.  A union bound proves (6.8).

It remains to simplify the global denominator.  For \(1\le x\le m_0\),

\[
 \binom{m_0}{x}
 =\prod_{j=0}^{x-1}\frac{m_0-j}{x-j}
 \ge\left(\frac{m_0}{x}\right)^x,
 \qquad
 x!\ge\left(\frac{x}{\mathrm{e}}\right)^x.
\]

The factorial inequality follows, for example, from
\(\sum_{j=1}^x\log j\ge\int_1^x\log t\,dt=x\log x-x+1\).
Consequently

\[
 (m_0)_x=x!\binom{m_0}{x}\ge(m_0/\mathrm{e})^x,          \tag{6.10}
\]

while \((m_0)_0=1=(m_0/\mathrm e)^0\) covers \(x=0\). Combining (6.10) with
\((d)_r\le d^r\) proves (6.9). Notice that (6.8) retains the single global falling
factorial before (6.9) is taken.
\end{proof}

\section{Exact partial diagonals}\label{exact-partial-diagonals}

\paragraph{Purpose of this section.}
The exact overlap representation of Section~6 still permits two partitions
to share whole classes. The definitions below assign a nonnegative marked
weight to each such common subprofile. We prove that the total marked
contribution is \(1+o(1)\), uniformly through the phase. This
partial-diagonal estimate is the input used in Section~8 to control the
endpoint-table sum. The counting identities in (7.1)--(7.4) are
finite-\(n\) identities; asymptotic estimates are explicitly labeled below.

A \emph{partial diagonal} is a specified collection of whole classes that
the two partitions have in common.  It is marked: one overlap may contain
several such collections and then contributes to every corresponding
nonnegative marked term.  This intentional overcount is what makes the
subsequent summation an upper-bound device.

For a selected common subprofile \(\ell=(\ell_i)\), put

\[
 \ell_\bullet=\sum_i\ell_i,\qquad
 m=\sum_i u_i\ell_i,\qquad
 b_\ell=\sum_i\binom{u_i}{2}\ell_i.                           \tag{7.1}
\]

Here \(0\le\ell_i\le k_i\) coordinatewise.  Since the full profile has
\(\sum_i u_i k_i=n\), this automatically gives \(m\le n\); thus every
factorial in the counting formula below has its ordinary combinatorial
meaning.

\paragraph{Exact marked decomposition.}
The signed first moment of such partial partitions occupying any
\(m\) vertices of \([n]\) is

\[
 Y_{\ell}^{\mathrm{sgn}}
 =2^{\ell_\bullet}\frac{n!}{(n-m)!\prod_i(u_i!)^{\ell_i}\ell_i!}
       2^{-b_\ell}.                                           \tag{7.2}
\]

For each type \(i\), choose the \(\ell_i\) marked slots in the first
partition and the \(\ell_i\) marked slots in the second, giving
\(\binom{k_i}{\ell_i}^2\) choices. The common blocks themselves have the
signed first-moment weight \(Y_\ell^{\mathrm{sgn}}\). In the normalized
second moment they are counted once rather than independently in each
marginal, so their contribution is the reciprocal of this weight.
Therefore the exact \emph{marked} common-subprofile weight is

\[
 D(\ell)=\frac{\prod_i\binom{k_i}{\ell_i}^2}
                 {Y_\ell^{\mathrm{sgn}}}.                         \tag{7.3}
\]

Different marked choices need not be disjoint: an overlap containing several
common blocks contributes to each corresponding marked choice. The sums below
are sums of these nonnegative marked weights, so no disjointness is used.
In particular,
\(D(0)=1\) and \(D(k)=1/\Exp{Z_k^{\mathrm{sgn}}}\).
To add one common block of type \(i\), first note that

\[
 \frac{\binom{k_i}{\ell_i+1}^2}{\binom{k_i}{\ell_i}^2}
 =\frac{(k_i-\ell_i)^2}{(\ell_i+1)^2},
\]

while (7.2) gives

\[
 \frac{Y_{\ell+e_i}^{\mathrm{sgn}}}
      {Y_\ell^{\mathrm{sgn}}}
 =\frac{2(n-m)_{u_i}}
        {u_i!(\ell_i+1)2^{\binom{u_i}{2}}}
 =\frac{2\mu_{u_i}(n-m)}{\ell_i+1}.
\]

Dividing these two ratios gives

\[
 \frac{D(\ell+e_i)}{D(\ell)}
 =\frac{(k_i-\ell_i)^2}
        {2(\ell_i+1)\mu_{u_i}(n-m)}.                     \tag{7.4}
\]

At the opposite endpoint, put \(h=k-\ell\),
\(v=\sum_i u_i h_i\), and let \(Z_h^{\mathrm{sgn}}(v)\) be the
signed first moment of a complete profile \(h\) on \(v\)
vertices. Exact factorial cancellation gives

\[
 D(k-h)=\frac{B(h)}{\Exp{Z_k^{\mathrm{sgn}}}},\qquad
 B(h)=\left(\prod_i\binom{k_i}{h_i}\right)Z_h^{\mathrm{sgn}}(v), \tag{7.5}
\]

and

\[
 \frac{\binom{k_i}{h_i+1}}{\binom{k_i}{h_i}}
 =\frac{k_i-h_i}{h_i+1},\qquad
 \frac{Z_{h+e_i}^{\mathrm{sgn}}(v+u_i)}
      {Z_h^{\mathrm{sgn}}(v)}
 =\frac{2\mu_{u_i}(v+u_i)}{h_i+1}.
\]

Therefore

\[
\frac{B(h+e_i)}{B(h)}
 =\frac{2(k_i-h_i)\mu_{u_i}(v+u_i)}{(h_i+1)^2}.          \tag{7.6}
\]

Equations (7.2)--(7.6) are exact finite-\(n\) counting identities.  We now
use them only through uniform asymptotic bounds on the three ranges of
selected mass.

In Heckel's admissible-phase argument, the corresponding control of partial
profiles enters through the tame-coloring estimates of
\citet[Lemmas~5.1--5.3, 6.3--6.5, and~7.20]{heckel-panagiotou-2023},
as used in \citet[Proposition~5]{heckel-2025-difference}; their hypotheses are the
source of the lower \(\mu_\alpha\)-window in that theorem.  The next
lemma is the replacement needed here.  It is proved directly for the
four-size signed profile, including both corners and every intermediate
mass, and no tame-profile theorem is invoked.

\begin{lemma}[All common subprofiles]
\label{lemma-7.1-all-common-subprofiles}

For the exact midpoint profile,

\[
 \sum_{0\le\ell_i\le k_i}D(\ell)=1+o(1).               \tag{7.7}
\]

\end{lemma}

We split the common-subprofile sum according to the vertex mass occupied by
the marked classes. The recurrence (7.4) controls the empty corner, the continuous rate
function controls the central range, and the reverse recurrence (7.6)
controls the full corner. For all sufficiently large \(n\), when
\(\eta<31/32\), the three ranges are the disjoint sets
\(m\le\eta n\); \(m>\eta n\) and \(n-m>n/32\); and
\(n-m\le n/32\). They are exhaustive, so no inclusion--exclusion or
boundary convention is hidden in the assembly.

\begin{proof}


\displayheading{Empty corner}

Let
\[
  \eta=\frac{\log\log n}{32\log n},
  \qquad
  \xi_i=\frac{k_i^2}{2\mu_{u_i}(n)},
  \qquad
  \Xi_{\mathrm{empty}}=\sum_i\xi_i.
  \tag{7.8}
\]
Equations (2.4), (2.8), and
$k_i=\Theta(n/\log n)$ imply, with one phase-independent constant,
\begin{equation*}
  \Xi_{\mathrm{empty}}
  =O\!\left((\log n)^{-(2/\log 2-1/2)}\right).
  \tag{7.9}
\end{equation*}
We now state explicitly how this activity at ambient size $n$ controls every
intermediate denominator in the selected-mass range.

Fix a final subprofile $\ell$ with selected mass $m\le\eta n$, and expose its
blocks in any order. At an intermediate selected mass $m'\le m$, the exact
ratio of first moments is
\[
  \frac{\mu_{u_i}(n)}{\mu_{u_i}(n-m')}
  =\frac{(n)_{u_i}}{(n-m')_{u_i}}
  =\prod_{r=0}^{u_i-1}
    \frac{n-r}{n-m'-r}.
\]
Uniformly in the phase, $u_i=O(\log n)$ whereas
$\eta n\gg\log n$. Hence, for all sufficiently large $n$,
$m'+r\le2\eta n$ throughout this product, and therefore
\[
  \frac{\mu_{u_i}(n)}{\mu_{u_i}(n-m')}
  \le(1-2\eta)^{-u_i}.
\]
The phase expansion and $\log 2>2/3$ give $u_i\le3\log n$ eventually. Since
$\eta\le1/4$ and $-\log(1-2\eta)\le4\eta$,
\begin{equation*}
  (1-2\eta)^{-u_i}
  \le\exp(4\eta u_i)
  \le\exp\!\left(\frac38\log\log n\right)
  =(\log n)^{3/8}.
  \tag{7.10}
\end{equation*}
Every threshold in this estimate is independent of the phase.

Suppose the current partial profile is $a$ and the next selected block has type
$i$. The exact recurrence (7.4), the inequality
$(k_i-a_i)^2\le k_i^2$, and (7.10) give
\[
  \frac{D(a+e_i)}{D(a)}
  \le
  \frac{(\log n)^{3/8}\xi_i}{a_i+1}.
\]
Multiplication along any exposure order gives the order-independent bound
\[
  D(\ell)
  \le
  \prod_i
  \frac{\bigl((\log n)^{3/8}\xi_i\bigr)^{\ell_i}}
       {\ell_i!}.
\]
All summands are nonnegative, so enlarging from the selected-mass range to the
whole nonnegative four-dimensional lattice and expanding the exponential
series yields
\begin{equation*}
  1
  \le
  \sum_{m\le\eta n}D(\ell)
  \le
  \exp\!\left((\log n)^{3/8}\Xi_{\mathrm{empty}}\right).
  \tag{7.11}
\end{equation*}
The lower bound is the empty marked subprofile.

For clarity, the exponent tends to zero with a fixed power margin. Indeed,
put $q=\log 2$. By (7.9),
\[
  (\log n)^{3/8}\Xi_{\mathrm{empty}}
  =O\!\left((\log n)^{-(2/q-7/8)}\right),
\]
and $q<7/10$ gives
\[
  \frac2q-\frac78
  >\frac{20}{7}-\frac78
  =\frac{111}{56}>0.
\]
Thus, for one deterministic sequence
$\varepsilon_n^{\mathrm{empty}}\to0$ independent of the phase,
\[
  1
  \le
  \sum_{m\le\eta n}D(\ell)
  \le e^{\varepsilon_n^{\mathrm{empty}}}
  =1+o(1).
\]



\displayheading{Central range}

Write
\[\begin{gathered}
p_i=k_i/k_{\mathrm{co}},\quad y_i=\ell_i/k_{\mathrm{co}},\quad
 Y=\sum_i y_i,\quad I=\sum_i i y_i,\\
z_i=p_i-y_i,\quad R=\sum_i z_i=1-Y,\quad
 T:=\sum_i i p_i=\alpha-\frac{n}{k_{\mathrm{co}}},\quad
 I_r=\sum_i i z_i=T-I.
\end{gathered}
\tag{7.12}\]
The residual vertex fraction is
\begin{equation*}
 \rho=\frac{n-m}{n}
 =R+\frac{I-TY}{\alpha-T}.
 \tag{7.13}
\end{equation*}

\displayheading{Uniform Stirling extraction}
Before approximating any factorial, equations (7.2) and (7.3) give the exact
identity
\[
\begin{split}
\log D(\ell)={}&
2\sum_i\{\log(k_i!)-\log(\ell_i!)-\log((k_i-\ell_i)!)\}\\
&-\ell_\bullet\log 2-\log(n!)+\log((n-m)!)
 +\sum_i\ell_i\log(u_i!)+\sum_i\log(\ell_i!)
 +(\log 2)b_\ell.
\end{split}
\]
We use the uniform form of Stirling's estimate
\[
  \log(t!)=t\log t-t+\sigma(t),
  \qquad
  |\sigma(t)|\le C_{\mathrm S}\log(t+2)
  \qquad(t\in\mathbb N_0),
\]
with the convention $0\log0=0$. Only a fixed number of factorials occur in
the preceding identity, and every argument is at most $n$. Therefore the total
remainder is $O(\!\log n)$ uniformly, including when some $\ell_i$ or
$k_i-\ell_i$ is zero. Using $n=k_{\mathrm{co}}(\alpha-T)$ gives
\begin{equation*}
\begin{split}
 \log D(\ell)={}&n\rho\log\rho\\
 &+k_{\mathrm{co}}\sum_i\{2p_i\log p_i
     -2(p_i-y_i)\log(p_i-y_i)
     -y_i\log y_i-y_i+y_iE_i\}
     +O(\!\log n),
\end{split}
 \tag{7.14}
\end{equation*}
where
\begin{equation*}
 E_i=\log k_{\mathrm{co}}+\log(u_i!)+u_i-u_i\log n
       +(\log 2)\binom{u_i}{2}-\log 2.
 \tag{7.15}
\end{equation*}

The entropy term in (7.14) is uniform even at the boundary of the coordinate
box. By (5.15) and the bounded tangent-rounding displacement, the exact integer
profile satisfies $p_i\ge c_p>0$ after decreasing $c_p$ once. Hence there is an
absolute $C_{\mathrm H}$ such that, for $c_p\le p\le1$ and $0\le y\le p$,
\begin{equation*}
  2p\log p-2(p-y)\log(p-y)-y\log y-y
  \le C_{\mathrm H}y\log\!\left(\frac{\mathrm e}{y}\right).
  \tag{7.14a}
\end{equation*}
Indeed, when $y\le p/2$, the mean-value theorem bounds
$p\log p-(p-y)\log(p-y)$ by $Cy$, while $-y\log y$ supplies the only
singular term. When $y\ge p/2$, one has $y\ge c_p/2$, and the assertion follows
by compactness. Finally,
\[
  \sum_i y_i\log\!\left(\frac{\mathrm e}{y_i}\right)
  \le
  Y\log\!\left(\frac{4\mathrm e}{Y}\right),
\]
by the entropy bound on four coordinates. This is
$O(Y\log(\mathrm e/Y))$ uniformly.

Since $u_{i+1}=u_i-1$, exact subtraction in (7.15) gives
\begin{equation*}
\begin{aligned}
 E_{i+1}-E_i
 &=\log n-\log(\alpha-i)-(\log 2)(\alpha-i)+\log 2-1\\
 &=-\frac{\log 2}{2}\alpha+O(1),
\end{aligned}
 \tag{7.16}
\end{equation*}
uniformly for $2\le i\le4$. For the second line, use
$(\log 2)\alpha/2=\log n-\log\log n+O(1)$ and
$\log(\alpha-i)=\log\log n+O(1)$.

Let $\bar E=\sum_i p_iE_i$. Applying Stirling once to the complete signed
first moment gives
\begin{equation*}
 -\frac1{k_{\mathrm{co}}}
   \log \Exp{Z_{\mathbf k}^{\mathrm{sgn}}}
 =\sum_i p_i\log p_i-1+\bar E+o(1).
 \tag{7.17}
\end{equation*}
The phase-uniform margin (5.19) implies $\bar E\le C_E$. To expose the affine
structure in (7.16), set
\[
  F_i:=E_i+\frac{(\log 2)\alpha}{2}i.
\]
Then $F_{i+1}-F_i=O(1)$. Because the support has only four consecutive
indices,
\[
  F_i=\sum_jp_jF_j+O(1)
     =\bar E+\frac{(\log 2)\alpha}{2}T+O(1).
\]
Consequently
\[
  E_i=\bar E-\frac{(\log 2)\alpha}{2}(i-T)+O(1),
\]
and therefore
\begin{equation*}
 \sum_i y_iE_i
 \le\frac{(\log 2)\alpha}{2}(TY-I)+O(Y).
 \tag{7.18}
\end{equation*}
Only the upper bound on $\bar E$ is used here.

It remains to compare the vertex fraction $\rho$ with the profile fraction
$R$. Since $2\le i,T\le5$,
\[
  |I-TY|=\left|\sum_i(i-T)y_i\right|\le3Y,
\]
so (7.13) gives
\[
  |\rho-R|\le\frac{3Y}{\alpha-T}=O(Y/\alpha).
\]
If $\rho\ge1/32$, then $Y\le1$ and the preceding bound implies
$R\ge1/64$ for all sufficiently large $n$. Hence $x\mapsto x\log x$ is
uniformly Lipschitz on the interval between $R$ and $\rho$. Moreover,
$-R\log R\le1-R=Y$. Thus
\[
\begin{aligned}
&\left|n\rho\log\rho
      -k_{\mathrm{co}}\alpha R\log R\right|\\
&\qquad\le
 k_{\mathrm{co}}(\alpha-T)C|\rho-R|
 +k_{\mathrm{co}}T|R\log R|
 \le Ck_{\mathrm{co}}Y,
\end{aligned}
\]
which proves
\begin{equation*}
 n\rho\log\rho
 =k_{\mathrm{co}}\alpha R\log R+O(k_{\mathrm{co}}Y).
 \tag{7.19}
\end{equation*}
Combining (7.14), (7.14a), (7.18), and (7.19) yields one absolute constant
$C$ such that
\begin{equation*}
 \log D(\ell)
 \le k_{\mathrm{co}}\alpha\Phi_T(z)
     +Ck_{\mathrm{co}}Y\log\!\left(\frac{\mathrm e}{Y}\right)
     +C\log n,
 \tag{7.20}
\end{equation*}
where
\begin{equation*}
 \Phi_T(z)=R\log R+\frac{\log 2}{2}(I_r-TR).
 \tag{7.21}
\end{equation*}
All constants and eventuality thresholds in this extraction are independent of
the phase.

\displayheading{Uniform rate negativity}
The four-deficit geometry gives
\begin{equation*}
 I_r-TR\le(5-T)R,
 \qquad
 I_r-TR=\sum_i(T-i)y_i\le(T-2)(1-R).
 \tag{7.22}
\end{equation*}
Put $q=\log 2$ and recall that $Y=1-R$. Multiplying the first inequality in
(7.22) by $Y$, the second by $R$, and adding gives
\begin{equation*}
 I_r-TR\le3RY.
 \tag{7.22a}
 \label{eq:partial-diagonal-combined-structural-v3}
\end{equation*}
For $0<R\le1$,
\[
 \log R\le\frac{2(R-1)}{R+1}.
\]
Indeed, for $x=(1-R)/R\ge0$, the function
$\log(1+x)-2x/(2+x)$ has derivative
$x^2/((1+x)(2+x)^2)\ge0$ and vanishes at zero. The positive expansion for
$\log 2$ gives
\[
  \frac23<q<\frac7{10}.
\]

If $1/64\le R\le3/4$, then
\[
\begin{aligned}
 \Phi_T
 &\le-\frac{2R}{1+R}Y+\frac{21}{20}RY\\
 &=-\left(\frac{2}{1+R}-\frac{21}{20}\right)RY
 \le-\frac{13}{8960}Y.
\end{aligned}
\]
Here $2/(1+R)\ge8/7$ and $R\ge1/64$. If $3/4\le R\le1$, then the phase
corridor gives $T<4$, so the second inequality in (7.22) implies
$I_r-TR\le2Y$. Since $\log R\le R-1=-Y$,
\[
  \Phi_T\le-RY+qY=(q-R)Y\le-\frac1{20}Y.
\]
Thus, throughout the scalar range,
\begin{equation*}
  \Phi_T(z)\le-\frac{Y}{5000}.
  \tag{7.23}
\end{equation*}
All rational inequalities in this two-range ledger are explicit; the
logarithmic inequality and four-deficit reduction are proved above.

\displayheading{Uniform summation over the central range}
Assume
\begin{equation*}
  m>\eta n,
  \qquad
  n-m>n/32.
  \tag{7.24}
\end{equation*}
Then $\rho>1/32$, and
\[
  \frac mn=\frac{\alpha Y-I}{\alpha-T}
  \le\frac{\alpha Y}{\alpha-T}
\]
implies
\[
  Y>\eta\frac{\alpha-T}{\alpha}\ge\frac\eta2
\]
for all sufficiently large $n$. Define the deterministic error ratio
\begin{equation*}
  \varepsilon_n^{\mathrm{diag}}
  :=
  \frac{C\log(2\mathrm e/\eta)}{\alpha}
  +
  \frac{2C\log n}{k_{\mathrm{co}}\alpha\eta}.
  \tag{7.24a}
\end{equation*}
Since
\[
  \alpha=\Theta(\log n),
  \qquad
  k_{\mathrm{co}}=\Theta(n/\log n),
  \qquad
  \eta=\Theta(\log\log n/\log n),
\]
one has $\varepsilon_n^{\mathrm{diag}}\to0$, uniformly in the phase. Because
$Y\ge\eta/2$, equations (7.20) and (7.23) give
\[
  \log D(\ell)
  \le
  -k_{\mathrm{co}}\alpha Y
   \left(\frac1{5000}-\varepsilon_n^{\mathrm{diag}}\right).
\]
For all sufficiently large $n$, the parenthesis is at least $1/10000$.
Moreover, $\alpha\eta=\Theta(\log\log n)$. Hence there is a phase-independent
$c>0$ such that
\begin{equation*}
  D(\ell)
  \le
  \exp\{-c k_{\mathrm{co}}\log\log n\}.
  \tag{7.25}
\end{equation*}
Define
\[
  \varepsilon_n^{\mathrm{central}}
  :=
  (k_{\mathrm{co}}+1)^4
  \exp\{-c k_{\mathrm{co}}\log\log n\}.
\]
There are at most $(k_{\mathrm{co}}+1)^4$ subprofiles, so the total
central-range contribution is at most
$\varepsilon_n^{\mathrm{central}}$. Since
$k_{\mathrm{co}}\log\log n\gg\log n$, this deterministic sequence tends to
zero, uniformly in the phase and with one eventuality threshold for the
complete phase.



\displayheading{Full corner}

Put $h=k-\ell$ and write
\[
  v(h):=\sum_i u_i h_i=n-m.
\]
The full corner is the range
\begin{equation*}
  v(h)\le\frac n{32}.
  \tag{7.26a}
\end{equation*}
We use the exact reverse recurrence (7.6) without replacing any ambient
falling factorial.

First make the uniform smallness of the residual first moments explicit. For
any $0\le w\le n$ and every one of the four block sizes,
\[
  \frac{\mu_{u_i}(w)}{\mu_{u_i}(n)}
  =\frac{(w)_{u_i}}{(n)_{u_i}}
  \le\left(\frac wn\right)^{u_i}.
\]
Equations (2.2) and (2.8) give
$\mu_{u_i}(n)\le n^{6+o(1)}$ uniformly in the phase. Moreover,
$u_i=(2/\log 2+o(1))\log n$, and
$\log 32=5\log 2$. Consequently, uniformly for $w\le n/32$,
\begin{equation*}
  \mu_{u_i}(w)
  \le
  n^{6+o(1)}32^{-u_i}
  =n^{-4+o(1)}.
  \tag{7.26}
\end{equation*}
In particular, there is one phase-independent threshold after which
\begin{equation*}
  \mu_{u_i}(w)\le n^{-3}
  \qquad
  (w\le n/32,\ 2\le i\le5).
  \tag{7.26b}
\end{equation*}
The relaxed exponent $-3$ leaves a fixed margin and is all that the recurrence
requires.

Fix a final residual profile $h$ satisfying (7.26a), and expose its blocks in
any order. Let $a$ be the current residual profile and
$v(a)=\sum_i u_i a_i$. If the next block has type $i$, then
$a_i<h_i$ and therefore
\[
  v(a)+u_i\le v(h)\le n/32.
\]
The exact recurrence (7.6), $k_i-a_i\le k_i\le n$, and (7.26b) now give
\[
  \frac{B(a+e_i)}{B(a)}
  =
  \frac{2(k_i-a_i)\mu_{u_i}(v(a)+u_i)}{(a_i+1)^2}
  \le2n^{-2}<1
\]
for all sufficiently large $n$. Hence $B$ never increases along the path from
$0$ to $h$, and
\begin{equation*}
  B(h)\le B(0)=1.
  \tag{7.27a}
\end{equation*}
By the exact complementary identity (7.5) and the phase-uniform first-moment
margin (5.19),
\begin{equation*}
  D(k-h)
  =\frac{B(h)}{\Exp{Z_k^{\mathrm{sgn}}}}
  \le e^{-c_Zk_{\mathrm{co}}}.
  \tag{7.27}
\end{equation*}
There are at most $(k_{\mathrm{co}}+1)^4$ residual profiles. Therefore
\[
  \sum_{v(h)\le n/32}D(k-h)
  \le
  \exp\{4\log(k_{\mathrm{co}}+1)-c_Zk_{\mathrm{co}}\}
  =o(1)
\]
uniformly in the phase. Equivalently, this sum is bounded by one deterministic
sequence $\varepsilon_n^{\mathrm{full}}\to0$.

\displayheading{Disjoint three-range assembly}

For sufficiently large $n$, one has $\eta<31/32$. The coordinate box of all
partial subprofiles is then the disjoint union of
\[
\begin{aligned}
  \mathcal E_n&:=\{\ell:m\le\eta n\},\\
  \mathcal C_n&:=\{\ell:m>\eta n,\ n-m>n/32\},\\
  \mathcal F_n&:=\{\ell:n-m\le n/32\}.
\end{aligned}
\]
Indeed, if $\ell\in\mathcal F_n$, then
$m\ge31n/32>\eta n$, so the full corner cannot meet the empty corner; after
excluding those two cases, the two strict inequalities defining
$\mathcal C_n$ are automatic. Thus no boundary term is counted twice and no
subprofile is omitted.

The three estimates proved above have the form
\[
\begin{aligned}
  1&\le\sum_{\ell\in\mathcal E_n}D(\ell)
      \le e^{\varepsilon_n^{\mathrm{empty}}},\\
  0&\le\sum_{\ell\in\mathcal C_n}D(\ell)
      \le\varepsilon_n^{\mathrm{central}},\\
  0&\le\sum_{\ell\in\mathcal F_n}D(\ell)
      \le\varepsilon_n^{\mathrm{full}},
\end{aligned}
\]
where every error sequence is deterministic, phase-independent, and tends to
zero. The lower bound in the first line is the term $D(0)=1$. Adding the three
disjoint sums gives
\[
  1
  \le
  \sum_{0\le\ell_i\le k_i}D(\ell)
  \le
  e^{\varepsilon_n^{\mathrm{empty}}}
  +\varepsilon_n^{\mathrm{central}}
  +\varepsilon_n^{\mathrm{full}}
  =1+o(1),
\]
which proves (7.7) with one eventuality threshold for the complete phase.
 \end{proof}

\section{Canonical high cells and endpoint-table comparison}
\label{sec:canonical-high-cells-v3}

This section estimates the contribution of the large overlap cells. The finite
combinatorics and the phase asymptotics are kept separate. We first sum the
set of all labeled stub matchings that realize a fixed high-cell demand. We
then compare
realized multiplicities with full containment, sum all positive deficits, and
regroup the full-containment references by endpoint table. Only in the final
subsection do we insert the four-size phase estimates. No residual local reward
or cycle-space factor is charged here; those factors belong to Section~9.

Section~7 has not subtracted or conditioned away overlaps containing common
whole classes. Instead, its weights provide a normalized one-sided reference
measure for the endpoint-table comparison below. Every overlap matrix remains present
and is represented exactly once by its canonical high-cell skeleton.

\subsection{All labeled realizations of a high skeleton}

Let $U$ be the largest class size in the fixed four-size profile and put
\[
  R_0:=\left\lfloor\frac U2\right\rfloor.
\]
A cell is \emph{high} if its multiplicity exceeds $R_0$. Two high cells cannot
share a row or a column: otherwise that row or column would contain more than
$U$ stubs. Hence all high cells form a canonical bipartite matching, denoted by
$P$.

For $e\in P$, let $s_e$ and $t_e$ be its endpoint block sizes and set
\[
  m_e:=\min\{s_e,t_e\},
  \qquad
  d_e:=|s_e-t_e|,
  \qquad
  h_e:=m_e-j_e,
\]
where $j_e$ is the realized multiplicity. Thus $m_e$ is the full-containment
multiplicity, $h_e$ is the deficit, and $j_e=m_e-h_e$. Since $j_e>U/2$ and
$m_e\le U$,
\begin{equation}
  2h_e<m_e.
  \label{eq:half-deficit-v3}
\end{equation}
Put $J:=\sum_{e\in P}j_e$.

For fixed \((P,j)\), let \(\mathfrak F(P,j)\) be the finite set of labeled
partial stub matchings that place exactly \(j_e\) pairs in every cell
\(e\in P\); no other stub pair is prescribed. We call
\(\mathfrak F(P,j)\) the \emph{physical fiber} of the high skeleton. Thus the
word ``fiber'' refers to the inverse image of \((P,j)\) under the map that
records only the selected cells and their multiplicities.

A size-$j$ partial matching between labeled blocks of sizes $s$ and $t$ can be
chosen in
\[
  \binom sj\binom tj j!
  =
  \frac{(s)_j(t)_j}{j!}
\]
ways: choose its two endpoint sets and then the bijection between them.

\begin{proposition}[Aggregate weight of a fixed high skeleton]
\label{prop:aggregate-high-weight-v3}
For a fixed matching support $P$ and multiplicity vector
$j=(j_e)_{e\in P}$, the sum of the bare weights over
\(\mathfrak F(P,j)\) is
\begin{equation}
  w(P,j)
  =
  \frac{
    \displaystyle\prod_{e\in P}(s_e)_{j_e}(t_e)_{j_e}
  }{
    \displaystyle(n)_J\prod_{e\in P}j_e!
  }
  \prod_{e\in P}g(j_e).
  \label{eq:aggregate-high-weight-v3}
\end{equation}
\end{proposition}

\begin{proof}
Because $P$ is a matching, distinct selected cells use disjoint row and column
stub sets. Therefore
\[
  |\mathfrak F(P,j)|
  =\prod_{e\in P}\frac{(s_e)_{j_e}(t_e)_{j_e}}{j_e!}.
\]
A prescribed set of $J$ stub pairs occurs with probability $(n)_J^{-1}$ in the
uniform bipartite configuration matching. Multiplying this probability by the
displayed cardinality and by the exposed reward
$\prod_{e\in P}g(j_e)$ gives \eqref{eq:aggregate-high-weight-v3}.
\end{proof}

From now on we write
\[
  w_{\mathrm{hi}}(P,j):=w(P,j).
\]
This is the same bare high-skeleton weight that appears in the exact
conditional decomposition of Section~9.

\begin{remark}
The matching hypothesis is load-bearing. If two selected cells shared a row or
a column, their endpoint choices would compete for the same stubs and the
one-cell fiber counts would not multiply.
\end{remark}

\subsection{Deficits and the single ambient loss}

For endpoint sizes $m$ and $m+d$, define
\begin{equation}
  R_{m,d}(h)
  :=
  \frac{\binom mh}
       {(d+1)(d+2)\cdots(d+h)}
  2^{-hm+h(h+1)/2},
  \qquad
  R_{m,d}(0):=1.
  \label{eq:local-deficit-ratio-v3}
\end{equation}
The product in the denominator is empty when $h=0$.

\begin{lemma}[Exact one-cell deficit ratio]
\label{lem:exact-local-ratio-v3}
For a selected high cell with $h\ge0$ and $2h<m$, the ratio of the aggregate
one-cell factor at multiplicity $m-h$ to its full-containment value at
multiplicity $m$ is $R_{m,d}(h)$.
\end{lemma}

\begin{proof}
The ratio of the physical matching counts is
\[
  \frac{(m)_{m-h}(m+d)_{m-h}}{(m-h)!}
  \frac{m!}{(m)_m(m+d)_m}
  =
  \frac{\binom mh}{(d+1)(d+2)\cdots(d+h)}.
\]
The local signed rewards satisfy
\[
  \frac{g(m-h)}{g(m)}
  =
  2^{-hm+h(h+1)/2}.
\]
Multiplying the two identities proves the claim.
\end{proof}

Write
\[
  w_{\mathrm{full}}(P):=w(P,(m_e)_{e\in P}),
  \qquad
  J_*:=\sum_{e\in P}m_e,
  \qquad
  H:=J_*-J=\sum_{e\in P}h_e.
\]
After the one-cell ratios have been extracted, the only factor that still
couples different cells is the ambient denominator:
\begin{equation}
  \frac{(n)_{J_*}}{(n)_J}
  =(n-J)_H
  \le n^H.
  \label{eq:one-global-falling-factorial-v3}
\end{equation}

\begin{lemma}[Aggregate deficit comparison]
\label{lem:aggregate-deficit-comparison-v3}
For every admissible deficit vector $h=(h_e)_{e\in P}$,
\begin{equation}
  w(P,m-h)
  \le
  w_{\mathrm{full}}(P)
  \prod_{e\in P}n^{h_e}R_{m_e,d_e}(h_e).
  \label{eq:aggregate-deficit-comparison-v3}
\end{equation}
\end{lemma}

\begin{proof}
Apply Lemma~\ref{lem:exact-local-ratio-v3} in each selected cell, apply
\eqref{eq:one-global-falling-factorial-v3} once, and use
$n^H=\prod_{e\in P}n^{h_e}$.
\end{proof}

The order of operations matters: first sum the physical fiber, and then apply
the single aggregate bound $(n-J)_H\le n^H$ from
\eqref{eq:one-global-falling-factorial-v3}.

\subsection{Summing all positive deficits}

For a selected cell $e$, enlarge the exact positive-deficit range to
\[
  A_e:=\{h\in\mathbb N:h\ge1,\ 2h<m_e\}.
\]
All weights are nonnegative, so the enlargement preserves the direction of the
upper bound.

\begin{lemma}[Optional-choice product]
\label{lem:optional-choice-product-v3}
Let $u_e(h)\ge0$ for $e\in P$ and $h\in A_e$. Then
\[
  \sum_{\omega}
    \prod_{e\in P}u_e(\omega_e)
  =
  \prod_{e\in P}
    \left(1+\sum_{h\in A_e}u_e(h)\right),
\]
where each $\omega_e$ is either the zero-deficit choice, of weight one, or one
positive deficit $h\in A_e$.
\end{lemma}

\begin{proof}
Expand the finite product. Each monomial records exactly one independent choice
in every selected cell, and every optional-choice vector occurs once.
\end{proof}

For $2h<m$, integrality gives $2h+1\le m$. Hence
\[
  \frac{h+1}{2}
  \le
  \frac{m+1}{4}
  \le
  m-\left\lfloor\frac{3m-1}{4}\right\rfloor.
\]
Multiplying by $h$ yields
\begin{equation}
  h\left\lfloor\frac{3m-1}{4}\right\rfloor
  \le
  hm-\frac{h(h+1)}2.
  \label{eq:three-quarter-budget-v3}
\end{equation}
Together with $\binom mh\le m^h$ and
\eqref{eq:local-deficit-ratio-v3}, this gives
\begin{equation}
  n^hR_{m,d}(h)
  \le
  \left(
    \frac{nm}{2^{\lfloor(3m-1)/4\rfloor}}
  \right)^h.
  \label{eq:three-quarter-charge-v3}
\end{equation}

For the remainder of this section, let $q=i+2$. The map
$i\mapsto q$ is a bijection from $\{0,1,2,3\}$ to the deficit labels
$\{2,3,4,5\}$ used in Sections~5--7, with inverse $q\mapsto q-2$. Define
\[
  \kappa_i:=k_{i+2},
  \qquad
  u_i:=\alpha-2-i,
  \qquad 0\le i\le3.
\]
Thus $\kappa_i$ is the number of endpoint blocks of size $u_i$. For
$i,j\in\{0,1,2,3\}$ let $m_{ij}:=\min\{u_i,u_j\}$ and put
\[
  \rho_{ij}
  :=
  \frac{n m_{ij}}{2^{\lfloor(3m_{ij}-1)/4\rfloor}},
  \qquad
  \rho_{16}:=\sum_{i,j=0}^{3}\rho_{ij}.
\]
Every local base in \eqref{eq:three-quarter-charge-v3} is at most
$\rho_{16}$.

\begin{lemma}[Coarse local deficit sum]
\label{lem:coarse-local-deficit-sum-v3}
Suppose $\rho_{16}\le1$. Then, for every selected cell $e$,
\[
  \sum_{h\in A_e}n^hR_{m_e,d_e}(h)
  \le
  (\alpha+1)\rho_{16}.
\]
\end{lemma}

\begin{proof}
If $e$ has endpoint type $(i,j)$, then
\eqref{eq:three-quarter-charge-v3} bounds its term of deficit $h$ by
$\rho_{ij}^h$. Since $h\ge1$ and
$\rho_{ij}\le\rho_{16}\le1$, this is at most $\rho_{16}$. Moreover,
$|A_e|\le m_e\le\alpha+1$. Summation proves the bound.
\end{proof}

\begin{proposition}[Fixed-support all-deficit bound]
\label{prop:fixed-support-all-deficit-v3}
Write $\mathbf m=(m_e)_{e\in P}$. If $\rho_{16}\le1$, then
\begin{equation}
  \sum_{\substack{h_e\in\mathbb N_0\ (e\in P)\\
                  m_e-h_e>R_0\ \text{for every }e}}
    w(P,\mathbf m-h)
  \le
  w_{\mathrm{full}}(P)
  \left(1+(\alpha+1)\rho_{16}\right)^{|P|}.
  \label{eq:fixed-support-all-deficit-v3}
\end{equation}
\end{proposition}

\begin{proof}
Use Lemma~\ref{lem:aggregate-deficit-comparison-v3}, enlarge the deficit range
to $A_e$, and apply Lemma~\ref{lem:optional-choice-product-v3} and
Lemma~\ref{lem:coarse-local-deficit-sum-v3}.
\end{proof}

Projection onto the row endpoint is injective on a block matching. Therefore,
if $K=\sum_{i=0}^{3}\kappa_i$ is the number of row-class slots, then
\begin{equation}
  |P|\le K.
  \label{eq:support-cardinality-v3}
\end{equation}

\subsection{Regrouping by endpoint table}

For $L=(\ell_{ij})\in\mathbb N^{4\times4}$, write
\[
  r_i=\sum_j\ell_{ij},
  \qquad
  c_j=\sum_i\ell_{ij}.
\]
Call $L$ \emph{feasible} if $r_i\le\kappa_i$ and $c_j\le\kappa_j$ for every
$i,j$. Such a table is realized by choosing the indicated row and column
block slots and pairing them cell by cell; conversely, every block matching
has a feasible table. For a feasible endpoint table $L$, define
\[
  W(L)
  :=
  \sum_{P:\,L(P)=L}w_{\mathrm{full}}(P).
\]
Equivalently, $W(L)$ is the sum of the full-containment atoms over all
block-slot matching supports with table $L$.

\begin{proposition}[Reference grouping]
\label{prop:reference-grouping-v3}
The full-containment references satisfy
\begin{equation}
  \sum_P w_{\mathrm{full}}(P)
  =
  \sum_{L\ \mathrm{feasible}}W(L).
  \label{eq:reference-grouping-v3}
\end{equation}
More generally, for every nonnegative function $F$ of the endpoint table,
\begin{equation}
  \sum_P w_{\mathrm{full}}(P)F(L(P))
  =
  \sum_{L\ \mathrm{feasible}}W(L)F(L).
  \label{eq:weighted-reference-grouping-v3}
\end{equation}
\end{proposition}

\begin{proof}
Partition the finite set of block-slot matching supports by the
value of $L(P)$. On each fiber the factor $F(L(P))$ is constant. Summing first
inside each fiber gives $W(L)$, and then summing over the feasible tables gives
the two identities.
\end{proof}

Whenever $\rho_{16}\le1$, combining
Proposition~\ref{prop:fixed-support-all-deficit-v3},
\eqref{eq:support-cardinality-v3}, and
\eqref{eq:reference-grouping-v3} yields
\begin{equation}
  \sum_{\text{high skeletons}}w_{\mathrm{hi}}
  \le
  \left(\sum_{L\ \mathrm{feasible}}W(L)\right)
  \left(1+(\alpha+1)\rho_{16}\right)^K.
  \label{eq:coarse-realized-table-reduction-v3}
\end{equation}

\begin{remark}[Sharper table-preserving form]
If every $\rho_{ij}\le1/2$, then
$\sum_{h\ge1}\rho_{ij}^h\le2\rho_{ij}$. Applying
\eqref{eq:weighted-reference-grouping-v3} with
$F(L)=\prod_{i,j}(1+2\rho_{ij})^{L_{ij}}$ gives
\[
  \sum_{\text{high skeletons}}w_{\mathrm{hi}}
  \le
  \sum_{L\ \mathrm{feasible}}
    W(L)\prod_{i,j}(1+2\rho_{ij})^{L_{ij}}.
\]
The proof below uses the coarser form
\eqref{eq:coarse-realized-table-reduction-v3}, which requires only the single
condition $\rho_{16}\le1$ and cleanly separates the deficit sum from endpoint
tables.
\end{remark}

\begin{remark}[Reusable finite core]
The aggregate-fiber formula, the single ambient-denominator comparison, the
optional-choice identity, and the endpoint-table regrouping are not specific
to four endpoint sizes. They apply to any finite endpoint alphabet for which
the selected cells form a matching. The four consecutive sizes enter only in
the numerical deficit charge and the endpoint-table estimate below.
\end{remark}

\subsection{Endpoint-table comparison}

Recall that $u_i=a-i$ and $\kappa_i=k_{i+2}$ for $0\le i\le3$, where
$a=\alpha-2$ and the corresponding deficit label is $i+2$.
For a full-containment table $L=(\ell_{ij})$, set
\[
  r_i=\sum_j\ell_{ij},
  \qquad
  c_j=\sum_i\ell_{ij},
  \qquad
  x_{ij}=\min\{u_i,u_j\},
  \qquad
  J(L)=\sum_{i,j}x_{ij}\ell_{ij}.
\]
Its exact reference weight is
\begin{equation}
  W(L)=
  \frac{\prod_i(\kappa_i)_{r_i}\prod_j(\kappa_j)_{c_j}}
       {\prod_{i,j}\ell_{ij}!}
  \frac1{(n)_{J(L)}}
  \prod_{i,j}
  \left[
    \frac{(u_i)_{x_{ij}}(u_j)_{x_{ij}}}{x_{ij}!}g(x_{ij})
  \right]^{\ell_{ij}}.
  \label{eq:endpoint-reference-weight-v3}
\end{equation}
For a one-sided margin vector $v$ satisfying $0\le v_i\le\kappa_i$, define
\begin{equation}
  D(v)=
  \frac{\prod_i(\kappa_i)_{v_i}^2}{\prod_i v_i!}
  \frac{\prod_i[u_i!g(u_i)]^{v_i}}{(n)_{m(v)}},
  \qquad
  m(v)=\sum_i u_iv_i.
  \label{eq:one-sided-diagonal-weight-v3}
\end{equation}
For comparison with Section~7, denote its common-subprofile vector by
$\ell^{\mathrm{pd}}=(\ell_q^{\mathrm{pd}})_{q=2}^{5}$. Under the bijection
$\ell_{i+2}^{\mathrm{pd}}=v_i$, the multiplicity and block size are
$k_{i+2}=\kappa_i$ and $\alpha-(i+2)=u_i$; hence (7.3) gives
$D(\ell^{\mathrm{pd}})=D(v)$. Below, $D(r)$ and $D(c)$ evaluate
\eqref{eq:one-sided-diagonal-weight-v3} on the margins in
$0\le v_i\le\kappa_i$.

For the endpoint pair $(i,j)$, put
\[
  s_{ij}:=\min\{u_i,u_j\},
  \qquad
  t_{ij}:=\max\{u_i,u_j\},
  \qquad
  d_{ij}:=t_{ij}-s_{ij},
\]
and define
\[
  Q_{ij}:=
  (n+1)^{d_{ij}/2}
  \frac{\sqrt{(t_{ij})_{d_{ij}}}}{d_{ij}!}
  2^{-\{d_{ij}s_{ij}+\binom{d_{ij}}2\}/2}.
\]
Thus $Q_{ii}=1$. Since $t_{ij}\le a$, $s_{ij}\ge a-3$, and
$d_{ij}\le3$,
\begin{equation}
  Q_{ij}\le\frac{(\tau_n^{\mathrm{end}})^{d_{ij}}}{d_{ij}!},
  \qquad
  \tau_n^{\mathrm{end}}=
  2^{3/2}\sqrt{\frac{(n+1)a}{2^a}}
  =O\!\left(\frac{(\log n)^{3/2}}{\sqrt n}\right).
  \label{eq:endpoint-eta-v3}
\end{equation}

Let
\[
  A_L=\frac{\prod_i r_i!}{\prod_{i,j}\ell_{ij}!},
  \qquad
  C_L=\frac{\prod_j c_j!}{\prod_{i,j}\ell_{ij}!},
  \qquad
  Q^L=\prod_{i,j}Q_{ij}^{\ell_{ij}}.
\]

\begin{lemma}[Endpoint-table product comparison]
For every feasible endpoint table $L$,
\begin{equation}
  W(L)^2
  \le
  \bigl(D(r)A_LQ^L\bigr)\bigl(D(c)C_LQ^L\bigr).
  \label{eq:endpoint-table-product-comparison-v3}
\end{equation}
\end{lemma}

\begin{proof}
The profile factors and the table factorials cancel exactly in the quotient of
$W(L)^2$ by $(D(r)A_L)(D(c)C_L)$. It remains to compare the ambient falling
factorials and the local full-containment atoms.

Fix one endpoint pair and abbreviate $s=s_{ij}$, $t=t_{ij}$,
$d=d_{ij}=t-s$, and $x_{ij}=s$. Its local atom in
\eqref{eq:endpoint-reference-weight-v3} is
\[
  \frac{(s)_s(t)_s}{s!}g(s)
  =
  \frac{t!}{d!}g(s).
\]
Consequently
\[
  \frac{
    \left[\frac{(s)_s(t)_s}{s!}g(s)\right]^2
  }{s!g(s)\,t!g(t)}
  =
  \frac{(t)_d}{(d!)^2}
  2^{-ds-\binom d2}
  =
  \frac{Q_{ij}^2}{(n+1)^d}.
\]
For all sufficiently large $n$, one has $s\ge\alpha-5\ge2$. Hence we may use
$g(t)/g(s)=2^{ds+\binom d2}$; the identity is also valid for $d=0$.

Moreover, $m(r),m(c)\ge J(L)$ and
\[
  m(r)+m(c)-2J(L)
  =
  \sum_{i,j}d_{ij}\ell_{ij}.
\]
Therefore
\[
  \frac{(n)_{m(r)}(n)_{m(c)}}{(n)_{J(L)}^2}
  \le
  (n+1)^{m(r)+m(c)-2J(L)}.
\]
Combining this inequality with the local ratio cell by cell gives
\[
\begin{aligned}
 \frac{W(L)^2}{(D(r)A_L)(D(c)C_L)}
 &\le
 (n+1)^{m(r)+m(c)-2J(L)}
 \prod_{i,j}
 \left(\frac{Q_{ij}^2}{(n+1)^{d_{ij}}}\right)^{\ell_{ij}}\\
 &=\prod_{i,j}Q_{ij}^{2\ell_{ij}}
  =(Q^L)^2,
\end{aligned}
\]
where the equality uses
\(m(r)+m(c)-2J(L)=\sum_{i,j}d_{ij}\ell_{ij}\). This is
\eqref{eq:endpoint-table-product-comparison-v3}.
\end{proof}

\begin{proposition}[Endpoint-table sum]
\label{prop:endpoint-table-sum-v3}
For all sufficiently large $n$, uniformly in the phase,
\begin{equation}
  \sum_{L\ \mathrm{feasible}}W(L)
  \le
  \exp\!\left\{O(\tau_n^{\mathrm{end}}K)\right\}
  \sum_rD(r)
  \le
  \exp\!\left\{O(\sqrt{n\log n})\right\}.
  \label{eq:endpoint-table-sum-v3}
\end{equation}
All sums in this proposition are finite.
\end{proposition}

\begin{proof}
By the arithmetic--geometric mean inequality and
\eqref{eq:endpoint-table-product-comparison-v3},
\[
  2W(L)
  \le
  D(r)A_LQ^L+D(c)C_LQ^L.
\]
Consider the first term. For a fixed row margin $r$, feasible tables form a
subset of all nonnegative integer tables with row margin $r$. Since
$A_LQ^L\ge0$, the multinomial theorem gives
\[
  \sum_{\substack{L\ \mathrm{feasible}\\
                  \operatorname{row}(L)=r}}A_LQ^L
  \le
  \sum_{\substack{L\in\mathbb N_0^{4\times4}\\
                  \operatorname{row}(L)=r}}A_LQ^L
  =
  \prod_i\left(\sum_jQ_{ij}\right)^{r_i}
  \le
  (1+C\tau_n^{\mathrm{end}})^K.
\]
After multiplication by $D(r)$ and summation over $r$, the first half is at
most $(1+C\tau_n^{\mathrm{end}})^K\sum_rD(r)$. The second half is identical after exchanging
rows and columns. Section~7 proves $\sum_rD(r)=1+o(1)$ uniformly in the phase.
This is a domination by a reference sum, not a deletion of common-class
overlaps. Finally, $K=\Theta(n/\log n)$ and
\eqref{eq:endpoint-eta-v3} imply
$\tau_n^{\mathrm{end}}K=O(\sqrt{n\log n})$.
\end{proof}

\subsection{Insertion of the phase estimates}

The phase satisfies
\[
  \alpha
  =
  \left(\frac{2}{\log 2}+o(1)\right)\log n.
\]
In particular, for all sufficiently large $n$,
\begin{equation}
  \frac52\log n\le\alpha.
  \label{eq:coarse-phase-corridor-v3}
\end{equation}
Every endpoint minimum satisfies $m_{ij}\ge\alpha-5$, and elementary floor
arithmetic gives
\[
  3\alpha-19
  \le
  4\left\lfloor\frac{3m_{ij}-1}{4}\right\rfloor.
\]
Using $\log 2>2/3$ and \eqref{eq:coarse-phase-corridor-v3}, we obtain
\begin{equation}
  \frac54\log n-\frac{19}{6}
  \le
  (\log 2)
  \left\lfloor\frac{3m_{ij}-1}{4}\right\rfloor.
  \label{eq:coarse-log-budget-v3}
\end{equation}
Hence
\[
  2^{\lfloor(3m_{ij}-1)/4\rfloor}
  \ge
  e^{-19/6}n^{5/4}.
\]
Since every $m_{ij}=O(\log n)$,
\begin{equation}
  \rho_{16}
  =O\!\left(\frac{\log n}{n^{1/4}}\right)
  \longrightarrow0.
  \label{eq:rho-sixteen-small-v3}
\end{equation}
Thus $\rho_{16}\le1$ for all sufficiently large $n$.

Define the bare high-skeleton sum by
\[
  \operatorname{BareSkeletonSum}_n
  :=
  \sum_{(P,j)\ \mathrm{feasible}}w_{\mathrm{hi}}(P,j),
\]
where the sum ranges over the canonical high-cell matching supports and their
admissible high multiplicities defined above.

\begin{proposition}[Bare high-skeleton estimate]
\label{prop:bare-high-skeleton-v3}
There is a deterministic sequence $\varepsilon_n\to0$ such that
\begin{equation}
  \operatorname{BareSkeletonSum}_n
  \le
  \exp\!\left\{
    \varepsilon_n\frac{n}{(\log n)^4}
  \right\}.
  \label{eq:bare-skeleton-final-v3}
\end{equation}
\end{proposition}

\begin{proof}
The deficit factor in \eqref{eq:coarse-realized-table-reduction-v3} satisfies
\[
  \log\left(1+(\alpha+1)\rho_{16}\right)^K
  \le
  K(\alpha+1)\rho_{16}
  =O(n^{3/4}\log n).
\]
Proposition~\ref{prop:endpoint-table-sum-v3} contributes
$O(\sqrt{n\log n})$ to the logarithm. Both quantities are
$o(n/(\log n)^4)$. Combining these estimates with
\eqref{eq:coarse-realized-table-reduction-v3} proves the proposition.
\end{proof}

No residual local reward or binary cycle-space factor has been included in
$\operatorname{BareSkeletonSum}_n$. Section~9 treats exactly those remaining
factors. All overlap matrices, including those containing whole common
classes, remain part of the exact conditional decomposition there.

\section{Residual attachments by matching restriction}
\label{sec:residual-attachments-v3}

Section~8 bounded the sum of the bare high-skeleton weights. It did not include
the rewards of unexposed cells or the binary cycle-space factor of the residual
support. We now condition on one feasible high skeleton and estimate precisely
those remaining factors. Section~7 has not removed any overlap from the exact
sum: its diagonal weights enter Section~8 only as a comparison measure for the
endpoint-table sum.

Throughout this section, write $Z:=Z_{\mathbf k}^{\mathrm{sgn}}$ for the
selected signed-profile witness count.

\subsection{Conditional decomposition}

Fix a canonical high skeleton with exposed block matching $M$, exposed
multiplicities $j$, and total exposed mass $J$. Put
\[
  m_0:=n-J.
\]
Let $(d_a)_a$ and $(d'_b)_b$ be the residual row and column degrees. Their two
sums equal $m_0$, and every degree is at most the phase cap $U$.

The unexposed pairs form a uniform bipartite configuration matching with these
residual degrees; when $m_0=0$, this means the unique empty matching. Let
$r'_{ab}$ be its cell counts and let $H_{\mathrm{res}}$ be the simple support
graph of cells with $r'_{ab}\ge2$. The residual event $\mathcal E(M,j)$ imposes
the cap $r'_{ab}\le\lfloor U/2\rfloor$ and forbids any further pair in a cell
of $M$. The latter no-return condition makes the canonical exposure unique.
More explicitly, for any fixed set $S$ of $J$ exposed stub pairs and any event
$\mathcal B$ of the residual matching,
\begin{equation}
  \mathbb P\{S\subseteq\Pi,\ \Pi\setminus S\in\mathcal B\}
  =
  \frac{1}{(n)_J}\,\mathbb P_{\mathrm{res}}(\mathcal B).
  \label{eq:exposed-residual-factorization-v3}
\end{equation}
After the physical fiber is summed, the residual attachment depends on the
block matching $M$ and the residual degree lists, but not on which labeled
stubs realize the exposed pairs. Define
\begin{equation}
  \mathcal A(M,j)
  :=
  \mathbb E_{\mathrm{res}}\!\left[
    \left(\prod_{a,b}g(r'_{ab})\right)
    2^{\beta(M\cup H_{\mathrm{res}})}
    \mathbf 1_{\mathcal E(M,j)}
  \right].
  \label{eq:residual-attachment-v3}
\end{equation}
Only the residual configuration matching is random in this expectation.

The exact overlap decomposition is
\begin{equation}
  \frac{\mathbb E Z^2}{(\mathbb E Z)^2}
  =
  \sum_{(M,j)}w_{\mathrm{hi}}(M,j)\,\mathcal A(M,j),
  \label{eq:exact-attachment-decomposition-v3}
\end{equation}
where the sum ranges over the finite family of feasible canonical high
skeletons. To justify the word \emph{exact}, start with an overlap matrix $r$.
Its cells above $\lfloor U/2\rfloor$ determine one and only one matching $M$
and multiplicity vector $j$. Removing the exposed pairs leaves one residual
matrix $r'$ satisfying $\mathcal E(M,j)$. Conversely, a feasible skeleton and
a residual matching satisfying $\mathcal E(M,j)$ reconstruct one overlap
matrix. The configuration-model cancellation proved in Section~8 factors its
probability and exposed local reward into $w_{\mathrm{hi}}(M,j)$ times the
residual law in \eqref{eq:residual-attachment-v3}; the remaining sign factor is
exactly $2^{\beta(M\cup H_{\mathrm{res}})}$. Thus no overlap is omitted and no
overlap is assigned to two skeletons.

The case $m_0=0$ is complete without introducing cell activities. The residual
matching is empty, every residual local factor is one, and
$H_{\mathrm{res}}=\varnothing$. Since $M$ is a matching,
$\beta(M)=0$, so
\begin{equation}
  \mathcal A(M,j)=1
  \qquad(m_0=0).
  \label{eq:zero-residual-attachment-v3}
\end{equation}
In all subsequent subsections we assume $m_0>0$. In view of
Proposition~\ref{prop:bare-high-skeleton-v3}, it remains to bound
$\mathcal A(M,j)$ uniformly in the positive-residual case.

\subsection{Threshold expansion and cell activities}

For every row--column pair $(a,b)$, put
\begin{equation}
  \theta_{ab}:=
  \frac{\mathrm e\,d_ad'_b}{m_0}.
  \label{eq:theta-v3}
\end{equation}
The intensity is defined on cells of $M$ as well as outside $M$; only the
activities will be suppressed on $M$. Let
\[
  \Delta_x:=g(x)-g(x-1),
  \qquad
  R:=\left\lfloor\frac U2\right\rfloor,
\]
and, for $(a,b)\notin M$, define
\begin{equation}
  \lambda_{ab}
  :=
  \sum_{x=3}^{R}\Delta_x\frac{\theta_{ab}^x}{x!},
  \qquad
  q_{ab}
  :=
  \frac{\theta_{ab}^2}{2}+\lambda_{ab}.
  \label{eq:lambda-q-v3}
\end{equation}
On $M$, set $\lambda_{ab}=q_{ab}=0$.

We first record the configuration estimate used in the expansion. For a finite
array of nonnegative demands $x=(x_{ab})$ supported outside $M$, write
$x_a=\sum_bx_{ab}$, $x'_b=\sum_ax_{ab}$, and $X=\sum_{a,b}x_{ab}$. Markov's
inequality applied to the product of binomial cell counts gives
\[
  \mathbb P_{\mathrm{res}}(r'_{ab}\ge x_{ab}\text{ for all }a,b)
  \le
  \frac{
    \prod_a(d_a)_{x_a}\prod_b(d'_b)_{x'_b}
  }{
    (m_0)_X\prod_{a,b}x_{ab}!}.
\]
If $X\le m_0$, then $(m_0)_X\ge(m_0/\mathrm e)^X$; if $X>m_0$, the event is
empty. Hence
\begin{equation}
  \mathbb P_{\mathrm{res}}(r'_{ab}\ge x_{ab}\text{ for all }a,b)
  \le
  \prod_{a,b}\frac{\theta_{ab}^{x_{ab}}}{x_{ab}!}.
  \label{eq:prescribed-cell-residual-v3}
\end{equation}
This joint estimate is the reason that cells sharing a row or a column may be
expanded simultaneously; no independence between cells is asserted.

For $0\le r\le R$,
\[
  g(r)=1+\sum_{x=3}^{R}\Delta_x\mathbf 1_{\{r\ge x\}}.
\]
Let $E_0$ be the set of all residual cell positions outside $M$, and let
$\mathcal C(M)$ be the family of even subsets of $M\cup E_0$. For
$F\in\mathcal C(M)$, define
\[
\begin{split}
  \Phi_F(r')
  :=
  &\prod_{e\in F\setminus M}
  \left(
    \mathbf 1_{\{r'_e\ge2\}}
    +
    \sum_{x=3}^{R}\Delta_x\mathbf 1_{\{r'_e\ge x\}}
  \right)\\
  &\times
  \prod_{e\in E_0\setminus F}
  \left(
    1+
    \sum_{x=3}^{R}\Delta_x\mathbf 1_{\{r'_e\ge x\}}
  \right).
\end{split}
\]
The first product contains the support threshold required when $e$ belongs to
the even set. On the event $\mathcal E(M,j)$, expansion of the cycle-space
cardinality and of all local rewards gives the exact identity
\[
  \left(\prod_{a,b}g(r'_{ab})\right)
  2^{\beta(M\cup H_{\mathrm{res}})}
  =
  \sum_{F\in\mathcal C(M)}\Phi_F(r').
\]

\begin{lemma}[Fixed even-set expansion]
\label{lem:fixed-even-set-expansion-v3}
For every $F\in\mathcal C(M)$,
\[
  \mathbb E_{\mathrm{res}}\!\left[
    \Phi_F(r')\mathbf 1_{\mathcal E(M,j)}
  \right]
  \le
  \prod_{e\in F\setminus M}q_e
  \prod_{e\in E_0\setminus F}(1+\lambda_e).
\]
\end{lemma}

\begin{proof}
Expand $\Phi_F$ into its nonnegative threshold monomials. In a cell of
$F\setminus M$, a monomial chooses either the threshold-two term or one higher
increment; it never chooses both. In a cell outside $F\cup M$, it chooses
nothing or one higher increment. Apply
\eqref{eq:prescribed-cell-residual-v3} once to the complete demand of each
monomial. The threshold-two contribution is $\theta_e^2/2$, the higher
increments sum to $\lambda_e$, and the empty choice contributes one. Every
monomial is nonnegative and $\mathbf 1_{\mathcal E(M,j)}\le1$, so omitting
the cap and no-return indicator can only increase the expectation.
\end{proof}

Summing over $F$ and then inserting the missing factors
$1+\lambda_e\ge1$ gives
\begin{equation}
  \mathcal A(M,j)
  \le
  \left(\prod_{e\in E_0}(1+\lambda_e)\right)
  \sum_{F\in\mathcal C(M)}
    \prod_{e\in F\setminus M}q_e.
  \label{eq:lambda-even-factorization-v3}
\end{equation}

\subsection{Restriction outside the exposed matching}

The next finite lemma is independent of random graphs.

\begin{lemma}[Restriction-product bound]
\label{lem:restriction-product-v3}
Let $E$ be a finite set, let $\mathcal C$ be a finite family of subsets of
$E$, and let $I\subseteq E$. Suppose that
\[
  A\longmapsto A\setminus I
\]
is injective on $\mathcal C$. For nonnegative activities $(q_e)_{e\in E}$,
\[
  \sum_{A\in\mathcal C}
    \prod_{e\in A\setminus I}q_e
  \le
  \prod_{e\in E\setminus I}(1+q_e).
\]
\end{lemma}

\begin{proof}
The restrictions $A\setminus I$ form a subfamily of the power set of
$E\setminus I$ and occur without repetition. Enlarge the sum to the full power
set and expand the finite product.
\end{proof}

Apply the lemma with $E=M\cup E_0$, $I=M$, and
$\mathcal C=\mathcal C(M)$. To verify injectivity, suppose that two even sets
have the same restriction outside $M$. Their
symmetric difference is an even subset of the matching $M$. Every nonempty
subset of a matching has a vertex of degree one, so this symmetric difference
must be empty. Therefore
\begin{equation}
  \sum_{F\in\mathcal C(M)}
    \prod_{e\in F\setminus M}q_e
  \le
  \prod_{e\in E_0}(1+q_e).
  \label{eq:even-family-product-v3}
\end{equation}

Since $0\le\lambda_e\le q_e$, equations
\eqref{eq:lambda-even-factorization-v3} and
\eqref{eq:even-family-product-v3}, together with $1+x\le e^x$, imply
\begin{equation}
  \mathcal A(M,j)
  \le
  \exp\!\left(2\sum_{e\in E_0}q_e\right).
  \label{eq:q-only-attachment-v3}
\end{equation}
Thus both the local rewards and the cycle-space cardinality are controlled by
one total residual activity.

\begin{remark}
Lemma~\ref{lem:restriction-product-v3} applies to any weighted set family whose
restriction map is injective. Here the family happens to be a binary cycle
space and injectivity follows solely from the fact that the deleted edge set is
a matching.
\end{remark}

\subsection{The intrinsic residual regime}

Assume
\begin{equation}
  2^U\le m_0^3.
  \label{eq:intrinsic-residual-regime-v3}
\end{equation}
Then $m_0\ge2^{U/3}$ and, by the degree cap,
\[
  \theta_{ab}
  \le
  \mathrm e U^2 2^{-U/3}.
\]

\begin{lemma}[Quadratic activity bound]
\label{lem:q-quadratic-v3}
There is an absolute constant $C_0$ such that, under
\eqref{eq:intrinsic-residual-regime-v3},
\begin{equation}
  q_{ab}\le C_0\theta_{ab}^2
  \label{eq:q-quadratic-v3}
\end{equation}
for every residual cell outside $M$.
\end{lemma}

\begin{proof}
Fix a cell and write $\theta=\theta_{ab}$. If $\theta=0$, then
$q_{ab}=0$, so the claim holds. Assume $\theta>0$. Since
$0\le\Delta_x\le g(x)$, it is enough to bound
\[
  a_x:=g(x)\frac{\theta^x}{x!},
  \qquad 3\le x\le R.
\]
For $x\ge3$,
\[
  \frac{a_{x+1}}{a_x}
  =
  \frac{2^x\theta}{x+1},
  \qquad
  \frac{a_{x+2}/a_{x+1}}{a_{x+1}/a_x}
  =
  2\frac{x+1}{x+2}>1.
\]
Thus $(a_x)$ is log-convex, and its maximum on $[3,R]$ occurs at an endpoint.
Consequently
\[
  \lambda_{ab}
  \le
  R(a_3+a_R).
\]
Now $a_3=(2/3)\theta^3$. Since
$R\theta=O(U^3 2^{-U/3})$, one has $Ra_3\le\theta^2$ for all sufficiently
large $U$.

For the other endpoint, using
$R=\lfloor U/2\rfloor$ and
$\theta\le\mathrm e U^2 2^{-U/3}$ gives
\[
  \log_2\theta
  \le
  \log_2\mathrm e+2\log_2U-\frac U3.
\]
For sufficiently large $U$, one has $R\ge3$ and $g(R)=2^{\binom R2-1}$. Hence
\[
\begin{aligned}
  \log_2\!\left(\frac{a_R}{\theta^2}\right)
  &=
  \binom R2-1+(R-2)\log_2\theta-\log_2(R!)\\
  &\le
  \left(\frac{U^2}{8}+O(U)\right)
  -\left(\frac{U^2}{6}+O(U)\right)
  +O(U\log U)\\
  &=
  -\frac{U^2}{24}+O(U\log U).
\end{aligned}
\]
The two quadratic terms come respectively from the local reward $g(R)$ and
from the factor $\theta^{R-2}$; the factorial term only improves the upper
bound. Hence $Ra_R\le\theta^2$ for all sufficiently large $U$. It follows that
$\lambda_{ab}\le2\theta^2$ in that range. For the finitely many remaining
values of $U$, the quotient $q_{ab}/\theta^2$ is bounded on the compact interval
$0\le\theta\le\mathrm e U^2 2^{-U/3}$, with limit $1/2$ at $\theta=0$.
Enlarging the constant proves the lemma.
\end{proof}

Because $\theta_{ab}$ was defined for every pair, the cell intensities satisfy
the exact identity
\[
  \sum_{a,b}\theta_{ab}^2
  =
  \frac{\mathrm e^2}{m_0^2}
  \left(\sum_a d_a^2\right)
  \left(\sum_b(d'_b)^2\right).
\]
Since every degree is at most $U$ and both degree sums equal $m_0$,
\[
  \sum_a d_a^2\le Um_0,
  \qquad
  \sum_b(d'_b)^2\le Um_0.
\]
On $M$ the activities are zero, while outside $M$ the preceding lemma applies.
Therefore
\begin{equation}
  \sum_{e\in E_0}q_e\le C_1U^2.
  \label{eq:total-q-v3}
\end{equation}
Together with \eqref{eq:q-only-attachment-v3}, this proves
\begin{equation}
  \mathcal A(M,j)\le\exp(CU^2)
  \qquad\text{if }2^U\le m_0^3.
  \label{eq:intrinsic-attachment-v3}
\end{equation}

\subsection{The complementary residual regime}

Suppose instead that $2^U>m_0^3$. Then
\begin{equation}
  m_0<2^{U/3}\le2^{\lceil U/3\rceil}.
  \label{eq:small-residual-mass-v3}
\end{equation}
We first bound the integrand pointwise. Since the integrand is nonnegative and
$0\le\mathbf 1_{\mathcal E(M,j)}\le1$, equation
\eqref{eq:residual-attachment-v3} remains bounded above after the indicator is
omitted. Each residual cell count is at most $U$, so
\[
  \prod_{a,b}g(r'_{ab})
  \le
  2^{\sum_{a,b}\binom{r'_{ab}}2}
  \le
  2^{(U-1)m_0/2}.
\]
Moreover, every edge of $H_{\mathrm{res}}$ uses at least two residual pairs,
so $|E(H_{\mathrm{res}})|\le m_0/2$. The restriction map from even subsets
of \(M\cup H_{\mathrm{res}}\) to subsets of \(E(H_{\mathrm{res}})\) is
injective by the matching argument in
Lemma~\ref{lem:restriction-product-v3}. Hence
\[
  2^{\beta(M\cup H_{\mathrm{res}})}
  \le
  2^{|E(H_{\mathrm{res}})|}
  \le
  2^{m_0/2}.
\]
Multiplication yields
\begin{equation}
  \mathcal A(M,j)
  \le
  2^{Um_0/2}
  \qquad\text{if }2^U>m_0^3.
  \label{eq:complementary-attachment-v3}
\end{equation}

\begin{proposition}[Uniform residual attachment]
\label{prop:uniform-residual-attachment-v3}
There is an absolute constant $C$ such that every feasible canonical high
skeleton satisfies
\[
  \mathcal A(M,j)
  \le
  \begin{cases}
    1,&m_0=0,\\
    \exp(CU^2),&m_0>0\text{ and }2^U\le m_0^3,\\
    2^{Um_0/2},&m_0>0\text{ and }2^U>m_0^3.
  \end{cases}
\]
\end{proposition}

\subsection{Explicit global logarithmic ledger}
\label{subsec:explicit-global-ledger-v3}

The preceding propositions already prove the required normalized second-moment
bound. We now record their error terms in a form that can be consumed directly
by the amplification theorem and by a theorem-facing formalization.

First combine the three deterministic partial-diagonal errors from Section~7:
\begin{equation*}
  \varepsilon_n^{\mathrm{pd}}
  :=
  e^{\varepsilon_n^{\mathrm{empty}}}-1
  +\varepsilon_n^{\mathrm{central}}
  +\varepsilon_n^{\mathrm{full}}.
  \tag{9.31}
\end{equation*}
Then $\varepsilon_n^{\mathrm{pd}}\to0$ uniformly in the phase, and the disjoint
three-range assembly gives
\begin{equation*}
  \sum_rD(r)\le1+\varepsilon_n^{\mathrm{pd}}.
  \tag{9.32}
\end{equation*}
Here and below every sum has the finite combinatorial domain specified in its
original definition.

Let
\[
  \tau_n^{\mathrm{end}}
  :=2^{3/2}\sqrt{\frac{(n+1)a}{2^a}},
  \qquad a=\alpha-2.
\]
This is the same $\tau_n^{\mathrm{end}}$ as in
\eqref{eq:endpoint-eta-v3}. For all sufficiently large $n$,
$\tau_n^{\mathrm{end}}\le1$. Since $Q_{ii}=1$ and each row has only three
off-diagonal positions, \eqref{eq:endpoint-eta-v3} gives the explicit row
bound
\begin{equation*}
  \sum_jQ_{ij}
  \le1+3\tau_n^{\mathrm{end}}.
  \tag{9.33}
\end{equation*}
Indeed, for $i\ne j$ one has
$Q_{ij}\le(\tau_n^{\mathrm{end}})^{d_{ij}}/d_{ij}!
\le\tau_n^{\mathrm{end}}$.

Repeating the two arithmetic--geometric-mean sums in the proof of
Proposition~\ref{prop:endpoint-table-sum-v3}, now with the explicit constant in
(9.33), gives
\begin{equation*}
  \sum_LW(L)
  \le
  (1+3\tau_n^{\mathrm{end}})^K
  (1+\varepsilon_n^{\mathrm{pd}}).
  \tag{9.34}
\end{equation*}
Combining this with the one-charge deficit reduction
\eqref{eq:coarse-realized-table-reduction-v3}, define
\begin{equation*}
\begin{split}
  \Gamma_n^{\mathrm{skel}}
  :={}&
  K\log\!\left(1+(\alpha+1)\rho_{16}\right)\\
  &+K\log\!\left(1+3\tau_n^{\mathrm{end}}\right)
  +\log\!\left(1+\varepsilon_n^{\mathrm{pd}}\right).
\end{split}
  \tag{9.35}
\end{equation*}
With the notation $M=P$, the skeleton sum in
\eqref{eq:exact-attachment-decomposition-v3} is precisely the
$\operatorname{BareSkeletonSum}_n$ defined in Section~8. For all sufficiently
large $n$,
\begin{equation*}
  \operatorname{BareSkeletonSum}_n
  \le e^{\Gamma_n^{\mathrm{skel}}}.
  \tag{9.36}
\end{equation*}
This inequality displays separately the unique ambient deficit charge, the
off-diagonal endpoint-table charge, and the partial-diagonal reference mass.

The phase estimates from Section~8 imply
\[
\begin{split}
  K\log\!\left(1+(\alpha+1)\rho_{16}\right)
  &\le K(\alpha+1)\rho_{16}
    =O(n^{3/4}\log n),\\
  K\log\!\left(1+3\tau_n^{\mathrm{end}}\right)
  &\le3K\tau_n^{\mathrm{end}}
    =O(\sqrt{n\log n}),\\
  \log(1+\varepsilon_n^{\mathrm{pd}})&=o(1).
\end{split}
\]
Consequently the deterministic normalized skeleton error
\begin{equation*}
  \varepsilon_n^{\mathrm{skel}}
  :=
  \frac{(\log n)^4}{n}\Gamma_n^{\mathrm{skel}}
  \longrightarrow0.
  \tag{9.37}
\end{equation*}

Let $C_{\mathrm{att}}$ be an absolute constant for which
\eqref{eq:intrinsic-attachment-v3} holds, and put $q:=\log 2$. Define
\begin{equation*}
  \Gamma_n^{\mathrm{att}}
  :=
  \max\!\left\{
    C_{\mathrm{att}}U^2,
    \frac q2U2^{U/3}
  \right\}.
  \tag{9.38}
\end{equation*}
Proposition~\ref{prop:uniform-residual-attachment-v3} implies the uniform
bound
\begin{equation*}
  \mathcal A(M,j)\le e^{\Gamma_n^{\mathrm{att}}}
  \tag{9.39}
\end{equation*}
for every feasible high skeleton. In the intrinsic regime this is precisely
\eqref{eq:intrinsic-attachment-v3}. In the complementary regime,
$m_0<2^{U/3}$ and
\[
  2^{Um_0/2}
  =\exp\!\left(\frac q2Um_0\right)
  \le\exp\!\left(\frac q2U2^{U/3}\right).
\]
The zero-residual case contributes one.

Since
\[
  U=O(\log n),
  \qquad
  2^U=\Theta\!\left(\frac{n^2}{(\log n)^2}\right),
\]
one has
\[
  \Gamma_n^{\mathrm{att}}
  =O\!\left(
    (\log n)^2+n^{2/3}(\log n)^{1/3}
  \right)
  =o\!\left(\frac{n}{(\log n)^4}\right).
\]
Thus the deterministic normalized attachment error
\begin{equation*}
  \varepsilon_n^{\mathrm{att}}
  :=
  \frac{(\log n)^4}{n}\Gamma_n^{\mathrm{att}}
  \longrightarrow0.
  \tag{9.40}
\end{equation*}

\begin{proposition}[Explicit normalized second-moment ledger]
\label{prop:explicit-normalized-second-moment-ledger-v3}
Define
\begin{equation*}
  \Lambda_n
  :=
  \Gamma_n^{\mathrm{skel}}+\Gamma_n^{\mathrm{att}}
  =
  \left(
    \varepsilon_n^{\mathrm{skel}}
    +\varepsilon_n^{\mathrm{att}}
  \right)
  \frac{n}{(\log n)^4}.
  \tag{9.41}
\end{equation*}
Then
\[
  \Lambda_n=o\!\left(\frac{n}{(\log n)^4}\right)
\]
and
\begin{equation*}
  1
  \le
  \frac{\mathbb E Z^2}{(\mathbb E Z)^2}
  \le e^{\Lambda_n}.
  \tag{9.42}
\end{equation*}
\end{proposition}

\begin{proof}
Insert (9.39) into the exact decomposition
\eqref{eq:exact-attachment-decomposition-v3}. Because all summands are
nonnegative, the uniform attachment factor may be taken outside the finite
skeleton sum. Equation (9.36) then gives
\[
  \frac{\mathbb E Z^2}{(\mathbb E Z)^2}
  \le
  e^{\Gamma_n^{\mathrm{att}}}
  \operatorname{BareSkeletonSum}_n
  \le
  e^{\Gamma_n^{\mathrm{att}}+\Gamma_n^{\mathrm{skel}}}
  =e^{\Lambda_n}.
\]
The lower bound is the nonnegativity-of-variance inequality. Equations (9.37)
and (9.40) give the stated order of $\Lambda_n$.
\end{proof}

Thus the global second-moment exponent is not an unspecified new error. It is
the sum of two deterministic ledgers with disjoint responsibilities:
\[
  \text{partial diagonals, deficits, and endpoint-table comparison}
  \quad+\quad
  \text{residual local rewards and cycle space}.
\]
No factor is charged in both ledgers.
\section{Rare-event amplification}\label{rare-event-amplification}

Proposition~\ref{prop:explicit-normalized-second-moment-ledger-v3} and (1.4)
give the seed

\[
 \Prob{Z_{\mathbf k}^{\mathrm{sgn}}>0}\ge e^{-\Lambda_n}.       \tag{10.1}
\]

A positive signed witness is a cocoloring with $k_{\mathrm{co}}$ classes, so
the seed implies

\[
 \Prob{\zeta(G_n)\le k_{\mathrm{co}}}\ge e^{-\Lambda_n}.        \tag{10.2}
\]

This event may still be rare. We turn it into a high-probability cocoloring
event by adding the quantity displayed in (10.5), which is later shown to be
$o(n/(\log n)^3)$ and hence smaller than the root separation. The
argument follows the seed-to-typical principle of
\citet[Theorem~1]{heckel-2025-difference}, using vertex-exposure concentration
as in \citet[Theorem~1]{scott-2008-2017}. The first lemma controls every
possible leftover vertex set simultaneously; the second amplifies an arbitrary
seed exponent $\Lambda_n$.

\begin{lemma}[Simultaneous leftover coloring]
\label{lemma-10.1-simultaneous-leftover-coloring}

There is an absolute \(C_0\) such that, with probability \(1-o(1)\), every
\(S\subseteq[n]\) satisfies

\[
 \chi(G_n[S])\le C_0\frac{|S|}{{\log n}}+n^{1/3}.               \tag{10.3}
\]

\end{lemma}

\begin{proof}
Let \(H\) be the complement of \(G_n\). Put
\(u_0=\lceil n^{1/4}\rceil\). For any fixed \(u_0\)-set, (1.5)
gives probability
\(\exp(-\Omega(u_0^2))\)
that its \(H\)-edge density is below \(1/4\). There are at
most \(\binom n{u_0}\) such sets, and for an absolute \(c>0\),

\[
 \binom n{u_0}e^{-cu_0^2}
 \le\exp\{u_0\log(\mathrm{e} n/u_0)-cu_0^2\}=o(1).
\]

Thus a union bound shows that, with
probability \(1-o(1)\), every \(u_0\)-set has density at
least \(1/4\). If \(S\) has size \(s\ge u_0\), double-counting each edge of
\(H[S]\) over the \(u_0\)-subsets containing it gives
\[
 \frac{1}{\binom{s}{u_0}}
 \sum_{\substack{T\subseteq S\\|T|=u_0}}
 \frac{e_H(T)}{\binom{u_0}{2}}
 =\frac{e_H(S)}{\binom{s}{2}}.
\]
Therefore every larger set also has density at least \(1/4\).

Let $S_0$ be any vertex set of size at least \(n^{1/3}\). While
$|S_t|\ge u_0$, choose $v_t\in S_t$ of maximum degree in $H[S_t]$, define
$S_{t+1}:=N_H(v_t)\cap S_t$, and put $s_t:=|S_t|$.
Density at least \(1/4\) implies
$|N_H(v_t)\cap S_t|\ge(s_t-1)/4$. Hence

\[
 s_{t+1}\ge(s_t-1)/4,\qquad s_t\ge4^{-t}s_0-1/3.         \tag{10.3a}
\]

The sets are nested and $S_{t+1}\subseteq N_H(v_t)$, so every later choice is
adjacent in $H$ to every earlier one. Thus the chosen vertices form a clique
in $H$. Since \(s_0\ge\lceil n^{1/3}\rceil\), for every
\(0\le t\le\lfloor {\log n}/(13\log4)\rfloor\), equation (10.3a) gives
\[
 s_t\ge n^{1/3-1/13}-\frac13
      =n^{10/39}-\frac13
      \ge n^{1/4}
\]
for all sufficiently large \(n\). Thus the procedure constructs
at least \(c\log n\) vertices for an absolute \(c>0\), and these vertices form
an independent set in \(G_n\).

For an arbitrary \(S\), repeatedly remove such independent sets and
give each a new color until fewer than \(n^{1/3}\) vertices
remain; color the rest singly.  Every removed set has at least
\(c\log n\) vertices, so the total number of colors is at most

\[
 \left\lceil\frac{|S|}{c\log n}\right\rceil+n^{1/3}.
\]

After enlarging the absolute constant $C_0$, this is (10.3), simultaneously for
every \(S\).
\end{proof}

\begin{lemma}[Amplification from a seed]
\label{lemma-10.2-amplification-from-a-seed}

There is an absolute constant $C>0$ and a deterministic
$\varepsilon_n=o(1)$ with the following property. Suppose deterministic
integers $k_n\ge0$ and deterministic reals $\Lambda_n\ge0$ satisfy

\[
 \Prob{\zeta(G_n)\le k_n}\ge e^{-\Lambda_n}.             \tag{10.4}
\]

The same $C$ and $\varepsilon_n$, independent of $k_n,\Lambda_n,r$, satisfy,
uniformly for every deterministic choice $r=r(n)>0$,

\[
\begin{split}
 \mathbb{P}\!\Bigg(\zeta(G_n)>k_n+C\bigg(
 &\frac{\sqrt{n\Lambda_n}+\sqrt{nr}}{\log n}
   +n^{1/3}+1\bigg)\Bigg)\\
 &\le e^{-r}+\varepsilon_n.                              \end{split}
\tag{10.5}
\]

\end{lemma}

\begin{proof}
Let $\mathcal E_n$ be the simultaneous event in
Lemma~\ref{lemma-10.1-simultaneous-leftover-coloring}, and set
$\varepsilon_n:=\mathbb P(\mathcal E_n^c)=o(1)$. This sequence is independent
of $k_n$, $\Lambda_n$, and $r$.

For an induced set \(W\), let
\(\zeta(W)=\zeta(G_n[W])\),
and define

\[
 S_k=\max\{|W|:\zeta(W)\le k_n\}.                       \tag{10.6}
\]

Expose the random graph in \(n-1\) independent vertex blocks, where
block \(v\) contains the edges from \(v\) to earlier vertices.
Changing one block changes \(S_k\) by at most one. Indeed, after deleting
the affected vertex, every feasible induced set loses at most one vertex,
and the two graph configurations agree on all remaining edges. Thus a
maximizer in either configuration yields a feasible set of size at least one
less in the other configuration. Therefore (1.3) applies.

Since \(S_k=n\) exactly when \(\zeta(G_n)\le k_n\),
(10.4) and the upper one-sided bounded-differences tail give

\[
 e^{-\Lambda_n}
 \le\Prob{S_k-\Exp{S_k}\ge n-\Exp{S_k}}
 \le\exp\!\left\{-\frac{2(n-\Exp{S_k})^2}{n-1}\right\}.
\]

Taking logarithms and rearranging gives

\[
 n-\Exp{S_k}\le\sqrt{(n-1)\Lambda_n/2}.                  \tag{10.7}
\]

The lower tail with radius \(\sqrt{(n-1)r/2}\) gives, outside an event of
probability at most \(e^{-r}\),

\[
 n-S_k\le
 \sqrt{(n-1)\Lambda_n/2}+\sqrt{(n-1)r/2}.                \tag{10.8}
\]

Choose a maximizing set \(W\) and put
\(V_{\mathrm{left}}=[n]\setminus W\). Concatenating a
\(k_n\)-part cocoloring of \(G_n[W]\) with an ordinary
coloring of \(G_n[V_{\mathrm{left}}]\) gives

\[
 \zeta(G_n)\le k_n+\chi(G_n[V_{\mathrm{left}}]).         \tag{10.9}
\]

On $\mathcal E_n$ and the event in (10.8), apply (10.3) to
$V_{\mathrm{left}}$. Since $|V_{\mathrm{left}}|=n-S_k$, substitution into
(10.9), followed by an enlargement of the absolute constant $C$, gives the
number of additional parts displayed in (10.5). A union bound gives the
failure probability $e^{-r}+\varepsilon_n$, proving (10.5) uniformly in the
three deterministic parameters.
\end{proof}

Fix the absolute constant $C$ and the sequence $\varepsilon_n$ from
Lemma~\ref{lemma-10.2-amplification-from-a-seed}, and apply that lemma to
(10.2). Put

\[
 r_n=\frac{\sqrt n}{(\log n)^2}.                              \tag{10.10}
\]

Proposition~\ref{prop:explicit-normalized-second-moment-ledger-v3} gives the
first estimate below. The definition of $r_n$ and elementary asymptotics give
the other two; in particular, $r_n\to\infty$:

\[
 \frac{\sqrt{n\Lambda_n}}{\log n}=o(n/(\log n)^3),\qquad
 \frac{\sqrt{nr_n}}{\log n}=o(n/(\log n)^3),\qquad
 n^{1/3}=o(n/(\log n)^3).                                       \tag{10.11}
\]

Define the deterministic sequence

\[
 a_n=C\left(
 \frac{\sqrt{n\Lambda_n}+\sqrt{nr_n}}{\log n}+n^{1/3}+1\right).
                                                               \tag{10.12}
\]

Then (10.11) and
Lemma~\ref{lemma-10.2-amplification-from-a-seed} give

\[
 a_n=o(n/(\log n)^3),\qquad
 \Prob{\zeta(G_n)>k_{\mathrm{co}}+a_n}
 \le e^{-r_n}+\varepsilon_n\longrightarrow0.            \tag{10.13}
\]

\section{Final assembly and the quantitative constant}
\label{sec:final-assembly-v3}

We now combine the chromatic lower location, the signed cocoloring upper
location, and the amplification estimate. The two final events are intersected
by a union bound; no independence is used.

\subsection{The phase-resolved gap}

Let $r_+(n)$ be the ordinary first-moment root and let
$r_4^{\mathrm{co}}(n)$ be the signed four-size root.
Section~\ref{the-four-size-signed-first-moment-advantage} gives
\begin{equation}
  r_+(n)-r_4^{\mathrm{co}}(n)
  =
  \left[
    \frac{(\log 2)^2}{4}A_4(\delta_n)+o(1)
  \right]
  \frac{n}{(\log n)^3},
  \qquad
  A_4(\delta):=\log 2-D_4(\delta).
  \label{eq:phase-resolved-root-gap-v3}
\end{equation}
The signed profile is placed at the rounded midpoint
\[
  k_{\mathrm{co}}
  =
  \left\lceil
    \frac{r_4^{\mathrm{co}}(n)+r_+(n)}2
  \right\rceil.
\]
At this fixed value of $k_{\mathrm{co}}$, the tangent correction in
Section~\ref{the-four-size-signed-first-moment-advantage} changes only the four
type multiplicities by $O(1)$. Here $i\in\{2,3,4,5\}$,
$u_i=\alpha-i$, and $k_i$ is the number of classes of size $u_i$. The
correction preserves exactly both
\[
  \sum_i k_i=k_{\mathrm{co}}
  \qquad\text{and}\qquad
  \sum_i u_i k_i=n.
\]
Thus there is no additional correction to the total number of classes. Hence
\begin{equation}
  r_+(n)-k_{\mathrm{co}}
  =
  \left[
    \frac{(\log 2)^2}{8}A_4(\delta_n)+o(1)
  \right]
  \frac{n}{(\log n)^3}.
  \label{eq:midpoint-retained-gap-v3}
\end{equation}
Indeed, (5.13) gives
$r_+(n)-k_{\mathrm{co}}
=\tfrac12(r_+(n)-r_4^{\mathrm{co}})+O(1)$.
The coefficient $1/4$ in
\eqref{eq:phase-resolved-root-gap-v3} is therefore multiplied by $1/2$, giving
$1/8$. The ceiling contributes the displayed $O(1)$; the tangent correction
preserves $k_{\mathrm{co}}$ itself.

Section~\ref{a-valid-unrestricted-lower-location-for-chi} constructs a
deterministic integer $k_\chi^-$ such that
\[
  \mathbb P\bigl(\chi(G_n)>k_\chi^-\bigr)\longrightarrow1
\]
and $k_\chi^-=r_+(n)+O(\log n)$. Combining this with
\eqref{eq:midpoint-retained-gap-v3} gives
\begin{equation}
  k_\chi^- - k_{\mathrm{co}}
  =
  \left[
    \frac{(\log 2)^2}{8}A_4(\delta_n)+o(1)
  \right]
  \frac{n}{(\log n)^3}.
  \label{eq:integer-location-gap-v3}
\end{equation}

The amplification result (10.13) gives a deterministic sequence
$a_n=o(n/(\log n)^3)$ such that
\[
  \mathbb P\bigl(\zeta(G_n)\le k_{\mathrm{co}}+a_n\bigr)
  \longrightarrow1.
\]
On this event and the chromatic lower event,
\[
  \chi(G_n)-\zeta(G_n)
  >k_\chi^- -k_{\mathrm{co}}-a_n.
\]
A union bound and \eqref{eq:integer-location-gap-v3} therefore give a
nonnegative deterministic sequence $\varepsilon_n^{\mathrm{gap}}\to0$ for
which
\begin{equation}
  \mathbb P\!\left(
    \chi(G_n)-\zeta(G_n)
    \ge
    \left[
      \frac{(\log 2)^2}{8}A_4(\delta_n)
      -\varepsilon_n^{\mathrm{gap}}
    \right]
    \frac{n}{(\log n)^3}
  \right)
  \longrightarrow1.
  \label{eq:phase-resolved-final-v3}
\end{equation}

Put $q:=\log 2$ and $\gamma_4:=\log(200/153)>0$. Lemma~\ref{lem:entropy-gap}
gives $A_4(\delta)>\gamma_4$ throughout the phase. Since
$\varepsilon_n^{\mathrm{gap}}\to0$, for all sufficiently large $n$,
\[
  \varepsilon_n^{\mathrm{gap}}<\frac{3q^2\gamma_4}{32}.
\]
Therefore the coefficient in \eqref{eq:phase-resolved-final-v3} is strictly
larger than
\[
  \frac{q^2\gamma_4}{8}-\frac{3q^2\gamma_4}{32}
  =\frac{q^2\gamma_4}{32}.
\]
It follows that
\[
  \mathbb P\!\left(
    \chi(G_n)-\zeta(G_n)
    \ge
    \frac{(\log 2)^2}{32}
    \log\!\left(\frac{200}{153}\right)
    \frac{n}{(\log n)^3}
  \right)
  \longrightarrow1.
\]
The displayed coefficient equals
\[
  \frac{(\log 2)^2}{32}
  \log\!\left(\frac{200}{153}\right)
  =0.004021983962242\ldots.
\]

This proves the main theorem. Since $n/(\log n)^3\to\infty$, the
chromatic--cochromatic difference tends to infinity with high probability
along the full sequence of integers.

\section*{Formal verification}
The exact theorem stated above is proved in Lean~4 by
\texttt{Erdos625.erdos625 : Erdos625Statement}. The formalization uses the
displayed constant
\((\log 2)^2\log(200/153)/32\), the full sequence of natural numbers, and the
finite uniform model of \(G(n,1/2)\). The modular project and a generated
single-file proof are publicly available at
\url{https://github.com/SamPetkov/Erdos/tree/main/625/formalization}; both use
the pinned Lean and Mathlib version \texttt{v4.31.0}. The final axiom audit reports only
\texttt{propext}, \texttt{Classical.choice}, and \texttt{Quot.sound}.

\renewcommand{\bibsection}{\section*{References}}
\begingroup
\small
\bibliographystyle{plainnat}
\bibliography{references}
\endgroup

\end{document}
